% 本文件由 LaTeX 排版 Agent 根据有效 translation.json 自动生成。
% author=Gregory of Nyssa; work_id=2907; title=论贞洁

\chapter{论贞洁}

% structure: p0001 source=h1 target=omitted reason=page-title-already-rendered-by-parent

% structure: p0002 source=h2 target=section reason=relative-heading-rank; base=h2

\section{第一章}

贞洁神圣的仪容确实珍贵，在所有以纯洁为美之准绳者看来；但唯独那些在争取这一高贵爱情对象的奋斗中，得天主恩宠眷顾和帮助的人，才拥有此仪容。它的赞美从与之相伴的名称中立刻就可听到；人们通常称它为\Quote{无玷}，这显示出其中的那种纯洁；因此我们可以通过这个对应词语来衡量这恩赐的高超，因为在众多德行的成果中，只有这一个享有了\Quote{无玷}之物的称号。倘若我们必须用赞辞来颂扬来自伟大天主的这一恩赐，祂的宗徒的话便足以赞美它；这些话虽简短，却使所有过分的颂词相形见绌；他只称那以这恩宠为装饰的人为\Quote{圣洁无瑕}。如今，如果这圣洁德行的成就就在于使人成为\Quote{无瑕而圣洁的}，并且这些称号以其最初和最强的意义被用来荣耀那不朽的天主，那么，还有什么比这更大的对贞洁的赞美呢？因为贞洁以某种方式\Emph{神化}那些分享其纯洁奥秘的人，使他们成为那唯独真实圣洁无瑕者的荣耀的分享者；他们的纯洁与不朽正是使他们与祂建立关系的方法。许多人撰写冗长的颂词和详细的论文，企图为这恩宠的奇妙增添点什么，依我看，他们无意中却适得其反；他们夸大其辞的铺陈方式反使其可信度受到质疑。大自然的伟大自有其令人惊叹的方式；它们不需要言辞的辩护：例如苍穹、太阳，或宇宙的任何其他奇观。在这下界的事务中，言辞无疑起到基础作用，赞美的技艺也确实赋予一种华丽的外表；以至于人类往往怀疑颂词所产生的惊叹不过是技巧的结果。因此，赞美贞洁的唯一充分方式将是表明这德行超越赞美，并通过我们的生活而非言辞来显明我们对它的钦佩。一个以此为主题欲求赞誉的人，似乎认为他自己的一滴汗水就能使浩瀚的海洋明显增大，如果他真的相信任何人类言辞都能为如此卓越的恩宠增添更多尊荣的话；那么他必然不是无知于自己的能力，就是无知于他所试图赞美的事物。\par

% structure: p0004 source=h2 target=section reason=relative-heading-rank; base=h2

\section{第二章}

要理解这恩宠的超凡卓越，确实需要深邃的思考。它包含在圣父不朽的观念中；这里一开始就有一个悖论，即在祂（圣父）内有圣子并毫无情欲地生了祂的童贞。它也包含在这独生圣子的性体内，祂首先奏响了这全部道德纯洁的基调；祂纯洁无欲的出生同样彰显出童贞。又是一个悖论：即圣子由童贞而为我们所知。它也见于圣神内在而不朽的纯洁；因为当你提及纯洁和不朽时，你便提到了童贞。它伴随整个超性存在；因其无欲，它常与天上的能力同在；从不与任何神圣事物分离，也绝不触及相反之物。凡其本能与意志都安于德行的人，都以此完美纯洁的不朽状态为美；凡属于相反品性的阶级，其现状与称呼都源于从纯洁的堕落。那么，什么样的言辞力量才能相称于这样的恩宠呢？怎能不担心任何人雄辩的努力会损害其内在价值的伟大呢？他所创造的对它的评价将低于听众先前的认识。因此，最好在此省略一切颂词；我们无法使言语达到主题的高度。相反，我们可以时常铭记天主的这份恩赐；我们的嘴唇可以常谈这祝福：尽管它是灵性存在和如此非凡卓越的属性，但因天主的爱，它被赐予了那些从肉情和血统中领受生命的人；当人性因情欲倾向而沉沦时，它伸出纯洁的援手，将其扶起并使之仰望高处。我想，正因如此，我们的主耶稣基督亲自——一切纯洁的泉源——并非通过婚姻来到世上。这是为了藉祂降生成人的方式揭示这伟大的奥秘：纯洁是神临在和来临的唯一完整标志，人若不完全远离肉体的情欲，便不能真正为自己获得纯洁。当那在基督内的天主性圆满透过无玷玛利亚闪耀时，发生在玛利亚身上的事，也同样发生在每个按规则度童贞生活的灵魂上。主确实不再以肉身临在；\BibleQuote{我们不再按血肉认识基督\BibleRef{格后 5:16}}；但祂以属灵的方式居住在我们内，并带着祂的父同来，正如福音在某处告诉我们的。那么，既然童贞具有如此重大的意义：它既在天上与众灵之父同在，在诸天能的行列中舞动，却又为人的救恩伸出双手；它既是将神性引下分享人性境遇的渠道，又为人的渴望添上翅膀飞向天上的事物，是神与人结合的纽带，通过它的中介使如此远隔的存在得以和谐——有什么言辞能强大到触及这奇妙的崇高呢？然而，如同无言语无感觉的受造物是可怕的；我们必须（若保持沉默）二选一：要么显得从未感受过童贞的特殊之美，要么显示自己固执地对一切美视而不见：因此我们同意按分配给我们这任务者的意愿，简要谈论这德行，在所有事上我们必须服从他。但不要指望我们有任何文采的展示；即使我们愿意，也许我们也不能产生，因为我们完全不善长那种写作。即使我们有这种能力，我们也不会以少数人的喜爱来换取多数人的造就。我认为，有见识的作家应该把首要目标定为不是超越所有其他作家而受赞赏，而是使自己和众多读者都得益处。\par

% structure: p0006 source=h2 target=section reason=relative-heading-rank; base=h2

\section{第三章}

但愿这次努力也能给我自己带来一些益处！如果我能够按照圣经分享耕田和打谷场上的果实，我会更乐意承担这劳苦；那么辛劳就会是乐趣。然而，现在我对童贞之美的认识在某种程度上对我是虚空无用的，正如那踹谷的牛嘴上的谷粒，或那从悬崖流下而渴者喝不到的水。那些仍然有能力选择更好的道路，并且没有被世俗生活的羁绊所阻碍的人，他们是幸福的，不像我们，现在被一道深渊与荣耀的童贞隔开：一旦踏上世俗生活，无人能攀登上去。我们只是他人幸福的旁观者，另一阶层幸福的见证人。即使我们想出一些关于童贞的恰当思想，我们也不会好于为富人餐桌提供甘美佳肴的厨师和帮厨，他们自己却没有份享用所预备的。如果情形不是这样，我们不必通过事后懊悔来认识善，那该多好啊！现在，\Emph{他们}是令人羡慕的，\Emph{他们}甚至超出了自己的祈祷和愿望，因为他们没有失去享受这些快乐的能力。我们就像那些以富有的社会对比自己的贫穷，因而更加烦恼和不满于现状的人。我们越准确地理解童贞的财富，就越要哀叹另一种生活；因为通过对比更好的事物，我们意识到它多么贫乏。我不仅指那些如此卓越生活的人在将来所存的赏报，也指他们活在此世时所拥有的赏报；因为如果有人决心准确衡量两种生活方式之间的差异，他会发现这差异几乎与天地之间的差异一样大。这句话的真实性可以通过观察实际事实来了解。\par

但在撰写这悲惨悲剧时，什么才是合适的开头？我们又如何真正将人生共有的种种邪恶呈现在眼前？人人都凭经验认识它们，但自然却以某种方式蒙蔽了实际受苦者，使他们心甘情愿地无视自己的处境。我们将从它最美好的甜蜜开始吗？那么，婚姻所期待的一切总和不就是获得愉快的伴侣吗？姑且假定这一点得到了；让我们勾勒一段各方面都最幸福的婚姻：出身显赫、财产充裕、年龄相配、正值青春鼎盛、深情厚爱、彼此都认为对方是最好的、那种渴望在爱上超越对方的甜蜜竞争；此外，还有名望、权力、广泛的声誉，以及其他一切。但请注意，即使在这福泽之下，也暗藏着不可避免的痛苦之火。我不说那种总是针对显贵者的嫉妒，以及笼罩在看似幸运者头上的易受攻击的威胁，还有那些不分享同样好运的人对优越者的天然憎恨，这些都使这些看似受宠的人度过一段焦虑的时光，其中痛苦多于快乐。我从这幅画面中略去这些，并假定对他们的嫉妒沉睡了；尽管要找到一个同时具备这两种福气——即超乎寻常的幸福和免于嫉妒——的生命并不容易。然而，如果可能，让我们假设一种免于所有这些考验的婚姻生活；看看那些拥有如此之多好运的人是否可能享受它。为什么？你会问，即使对他们幸福的嫉妒也触及不到他们时，还有什么烦恼留下？我断言，正是这环绕他们生活的甜蜜本身，就是点燃痛苦的火花。他们始终是人，脆弱而必朽；他们不得不看着祖先的坟墓；因此，痛苦与他们的存在密不可分，只要他们有丝毫反省的能力。这种对死亡的持续期待，虽无确凿征兆可确定，却以对未来可怕的未知悬在他们头上，扰乱了他们当下的喜悦，用对未来的恐惧使之蒙上阴影。如果能在经验到来之前就学到经验的结果，或者一旦踏上这条路，就有可能通过其他推测方法审视现实，那么，会有多少人从婚姻逃向贞女生活啊！他们会多么小心和急切地避免陷入那张束缚人的网罗，在那里，没有人确切知道网是如何刺痛人的，直到他们实际进入！如果你能毫无危险地做到这一点，你会看到那里有许多对立面结合在一起：笑容融化成眼泪，痛苦与快乐交织，死亡总是悬在出生的孩子头上，用手指点在每一个最甜美的快乐上。每当丈夫看着心爱的面容时，分离的恐惧便紧随那目光。若他聆听甜美的声音，心中便想到某天他将听不到它。当他因注视她的美貌而喜悦时，他最为因预感哀悼她而战栗。当他注意到所有那些对青年如此珍贵、轻率者所追求的魅力——眼睑下的明亮眼睛、弯弯的眉毛、带着甜美酒窝微笑的面颊、嘴唇上绽放的自然红润、金发缠绕在头上熠熠生辉，以及所有那转瞬即逝的优雅——那么，即使他不太反省，他也必定在内心深处有一个念头：所有这些美丽终将消融，化为乌有，在这番炫耀之后变成可憎而难看的骨头，不留痕迹、纪念或残余的生机。当他想到这些时，他能快乐地生活吗？他能信任这些仿佛永远属于他的宝藏吗？不，显然他会像被梦境愚弄般摇摆不定，对生命的信心动摇，不再把所见视为己有。如果你对事物本相有全面的了解，就会明白此生中没有一样东西是它看起来的那样。它通过我们想象的幻象向我们展示一个事物，而非另一个。人们张口凝视着它，它用希望嘲弄他们；它暂时隐藏在这虚幻的表象之下；然后，在生命的逆转中，它突然显露出与人们愚蠢心中因欺诈而产生的希望所描绘的完全不同。对于那些反省此事的人而言，生命的甜蜜似乎值得享受吗？他还能真正感受到它，以便对所拥有的财富感到喜悦吗？他是否会因不断担忧某种逆转而心烦意乱，虽有使用却无享受？我只提一下那些因愚蠢思考习惯而被注意到的预兆、梦境、征兆以及类似之物，它们总使人害怕最坏的情况。但年轻妻子的分娩时刻来临；这情况不被视为孩子出生，而被视为死亡的临近；在生育中，人们预期她会死；而且，这悲伤的预感常常成真，在设宴庆生之前，在尝到任何期待的喜悦之前，他们就得将欢乐变为哀悼。仍在爱情的热病中，仍在深情之爱的顶峰，尚未把握生命中最甜蜜的礼物，仿佛在梦中景象中，他们突然从所拥有的一切中被夺走。但接下来呢？家仆如同胜利的敌人，拆毁洞房；他们为葬礼布置它，但它现在是死亡之房；他们徒劳地哀哭捶胸。然后是往日的回忆，诅咒那些劝婚的人，指责那些没有阻止的亲友，将过错推给父母（无论他们是生是死），对人性命运的苦涩爆发，对整个自然进程的控诉，甚至对天主眷顾的抱怨和指控；人内心的战争，与那些规劝者争斗；对最骇人的言语和行为毫无反感。有些人这种心态持续，他们的理智更完全被悲伤吞噬；\Emph{他们的}悲剧结局更悲惨，受害者无法忍受活过这场灾难。\par

与其如此，不如让我们设想一个较为幸福的状况。分娩的危险已经过去；他们生了一个孩子，恰是父母美丽的肖像。难道忧愁的机会因此而减少了吗？反而增加了；因为父母保留着之前所有的恐惧，此外又增添了为孩子的担忧，惟恐在抚养过程中发生什么意外，例如一次严重的事故，或因某种不幸而染上疾病、发热，或任何危险的病症。父母双方都同样分担这些；但谁能述说妻子特有的焦虑呢？我们省略了最显而易见、人人都能理解的，如怀孕的疲惫、分娩的危险、哺乳的操劳、为子女而心碎，以及如果她是多子女的母亲，她的灵魂被分割成与子女数目相等的碎片；她对发生在子女身上的一切感受的温柔。这是每个人都很了解的。然而天主的圣言告诉我们，她不是自己的主人，只在她婚姻所立为主的那人身上找到她的资源；因此，即使她只被单独留下片刻，她就感到好像与自己的头分离了，难以忍受；她甚至将丈夫的短暂缺席视为她寡居的前奏；她的恐惧使她立刻放弃一切希望；于是，她那充满惊恐悬念的眼睛总是盯着门口；她的耳朵总是忙着倾听别人的耳语；她的心被恐惧刺痛，即使尚未传来任何坏消息，也几乎要爆裂；门口的一点声响，无论是想象还是真实的，都成了恶讯的信使，突然震撼她的灵魂；很可能外面一切平安，毫无可惧的原因；但她那昏厥的心灵比任何消息都快，将她的幻想从好消息转向绝望。因此，即使是最受眷顾的人也是如此生活，他们并非完全值得羡慕；他们的生活无法与童贞的自由相比。然而这仓促的描绘省略了许多更令人痛苦的细节。常常这位少妇，刚结婚，依然闪耀着新娘的优雅，也许当新郎进入时还会脸红，羞怯地偷偷瞥他一眼，而此时情欲因羞耻而无法表露，更加炽烈，却突然不得不承受可怜孤独寡妇的名号，被称为一切可悲的事物。死亡瞬间降临，将那位穿着白色华服的光彩造物变为黑衣哀悼者。它卸下新娘的饰物，以丧服的颜色为她穿衣。昔日欢快的房间陷入黑暗，侍女们唱起长长的挽歌。她憎恨那些试图减轻她悲伤的朋友；她不肯进食，日渐消瘦，在灵魂深深的沮丧中，只强烈渴望着死亡，这渴望常常持续到死亡来临。即使假设时间结束了这悲伤，但另一个悲伤还会来临，无论她是否有孩子。如果有，孩子便成了孤儿，作为怜悯的对象，重新唤起她失落的记忆。如果没有孩子，那么她失去的丈夫的名字便被连根拔起，这悲伤比那表面上的安慰更大。关于寡居的其他特殊悲伤，我将不多说；因为谁能精确地一一列举呢？她在亲属中发现敌人。有些人甚至利用她的痛苦。另一些人则对她的损失幸灾乐祸，带着恶意的喜悦看着家庭破碎、仆人的无礼以及在这种情况下可见的其他苦难，这些苦难比比皆是。因此，许多妇女被迫再次冒险尝试同样的事情，因为无法忍受这痛苦的讥讽。仿佛她们能通过增加自己的痛苦来报复侮辱！另一些妇女，记着过去，宁愿忍受一切，也不愿再次陷入同样的麻烦。如果你想知道这婚姻生活的一切考验，请听那些真正了解它的妇女们。她们如何祝贺那些从一开始就选择了童贞生活的人，不必通过经验得知更好的道路，即童贞可以抵御所有这些祸患，没有孤儿状态，没有寡居可哀；它常在与不朽新郎的同在中；它总有虔诚的后嗣可欢喜；它不断看见一个真正属于自己的家，充满各种珍宝，因为主人常住在那里；在这种情况下，死亡不会带来分离，而是与所渴望的那一位结合；因为当（灵魂）离去时，它就与基督同在，正如宗徒所说。\newline{} 但现在是时候了，既然我们已经一方面审视了那些命运幸福者的感受，就要揭示另一种生活，在那里贫穷、逆境以及人类必须忍受的一切其他邪恶都是固定的状态；我指的是畸形、疾病以及所有终身的痛苦。那生活在自己之内的人，要么完全避开它们，要么能轻易承受，拥有专注于自身而不被分心的心灵；而将自我分给妻子和孩子的人，往往连对自己状况的懊悔都无暇顾及，因为对亲人的忧虑充满他的心。但重复每个人都知道的事是多余的。如果连看似繁荣的事物都捆绑着如此痛苦和厌倦，那么对相反的情况我们还能期待什么呢？任何试图将其呈现于我们眼前的描述都将不及现实。但也许我们可以用很少的几句话来宣告其悲惨的深度。那些命运与所谓繁荣相反的人，其悲伤也来自与繁荣相反的原因。繁荣的生活被对死亡的期待或死亡本身所破坏；但这些人的悲惨在于死亡拖延了它的到来。因此，这些生活被相反的感受广泛分隔；虽然同样没有希望，却归于同一个结局。婚姻赋予世界的恶果是如此多面、如此奇特不同。痛苦总是存在，无论孩子出生，还是永远不能期待，无论他们活着还是死去。有人孩子众多却没有足够的生活资料；有人辛勤劳作积攒了巨大财富却没有继承人，便将别人的不幸视为福气；事实上，每个人都渴望自己看见别人悔恨的东西。再者，有人因死亡失去心爱的儿子；另有一个儿子活着却是败类；两者同样可怜，尽管一个为儿子的死哀伤，另一个为儿子的生哀伤。我也不再多说家庭嫉妒和争吵如何以可悲和灾难性的方式结束，这些争吵源于真实或想象的原因。谁能完全详述所有这些细节呢？如果你想了解人类生活是如何被这些邪恶缠绕，你不必追溯到那些为戏剧诗人提供主题的古老故事。\Emph{它们}因其骇人听闻的过分而被视为神话；其中有谋杀和吃子、杀夫、杀母杀兄、乱伦结合，以及各种天性的扰乱；然而古老的编年史家讲述以这些恐怖告终的故事时，是从婚姻开始的。但撇开这一切，只注视在这人生舞台上上演的悲剧；是婚姻为那里提供了演员。去法庭，通读那里的法律；然后你就会知道婚姻的羞耻秘密。正如当你听到医生解释各种疾病时，你通过了解身体易患的痛苦数量和种类而理解其苦难，同样，当你细读法律，读到其中与对应刑罚相连的婚姻中奇怪的犯罪种类时，你将对其性质有一个相当准确的认识；因为法律不为不存在的邪恶提供补救，正如医生不为从未知道的疾病提供治疗一样。\par

% structure: p0010 source=h2 target=section reason=relative-heading-rank; base=h2

\section{第四章}

但我們無需再用這狹窄的方式展示這種生活的弊端，彷彿其禍患的數量僅限於姦淫、紛爭和陰謀。我認為，我們應當採取更高、更真實的觀點，直接說：生活中所有邪惡，若人甘心被這條道路的鎖鏈束縛，那麼任何邪惡都無法掌控他。我們所說的話的真實性將如此顯明。一個以潔淨的心靈之眼看穿幻象，將自己提升於掙扎的世界之上，並用宗徒的話語，輕視這一切為糞土，因戒絕婚姻而在某種程度上自我放逐於人類生活之外的人——他與人類的罪惡毫無干係，如貪婪、嫉妒、忿怒、仇恨及諸如此類。他脫離了這一切，完全自由平安；沒有什麼會激起鄰人的嫉妒，因為他不抓取那些世間嫉妒所圍繞的對象。他已將自己的生命提升到世界之上，將美德視為唯一寶貴的財富，在無痛的平安與寧靜中度過日子。因為美德是一種財富，即使所有人都按各自的能力分享它，它對渴求者而言也總是豐盈有餘；不像這地上的土地佔有，人們將其劃分，一方增加則另一方減少，以致一人的所得便是另一人的損失；由此產生爭奪大份的鬥爭，源自人們憎恨受騙。但\Emph{這}財富的更大擁有者從不被嫉妒；搶奪大份的人並不損害要求平等分享的人；按各自的能力，每個人都能滿足這高貴的渴望，而那些已擁有者的美德財富也不會耗盡。因此，一個人若只注目於這樣的生活，以人無法設限的美德為唯一寶藏，他絕不會容忍自己的靈魂屈從於眾人所走的任何低賤途徑。他不會羨慕地上的財富、人的權力，或愚昧所追求的任何事物。的確，若他的心思仍如此低下，他就不是我們初學者的行列中的一員，我們的話語也不適用於他。但若他的思念在上，彷彿與天主同行，他將被提升脫離這一切錯誤的迷宮；因為導致這一切的原因——婚姻——未曾觸及他。如今，欲居人之上是驕傲的大罪，說這是罪惡荊棘的一切根源之根，也離事實不遠；但驕傲大多始於婚姻相關的理由。為表明我的意思，我們通常發現貪婪的人將過錯推給最親近的親屬；狂熱追求名聲和野心的人通常讓家庭為此罪負責：\Quote{他不可被視為不如他的祖先；他必須為後代留下自己的歷史記錄，從而被後世視為偉大}。同樣，靈魂的其他疾病——嫉妒、惡意、仇恨等——也與此原因相關；它們存在於那些熱衷於此生事物的人中。逃離此世的人，猶如從高高的瞭望塔俯瞰人類的景象，憐憫這些虛榮的奴隸，因他們盲目地看重身體的福祉。他看到某個顯要人物因公眾榮譽、財富和權力而自負，只嘲笑如此膨脹的愚昧。他依照聖詠作者的計算，給人類生命以最長的年數，然後將這原子般的瞬間與無盡的世代相比，憐憫那些為如此低賤、短暫、易朽的事物而激動的人的虛浮榮耀。在這一切事物中，有什麼值得眾人所追求的可羨慕之物——榮譽？得勝者得到了什麼？無論受不受榮譽，凡人仍是凡人。擁有許多土地的人最終得到什麼好處？除非愚昧人以為那從不屬於他的東西是他自己的，似乎在他的貪婪中不知\Quote{大地屬於上主，及其中的萬物}，因為\Quote{天主是全地的君王}。正是擁有的激情給人以虛假的主權去支配那永遠不屬於他們的東西。\Quote{大地}，明智的訓道者說，\BibleQuote{永遠存在}\BibleRef{訓 1:4}，服務於每一代人，先後生於其上的世代；但人雖然連自己的主人都算不上，由造物主的旨意不知不覺地帶入生命，又在不情願時被帶走，卻在極度虛榮中自以為是它的主人；他們時而生、時而死，卻統治那永續存在之物。反思這一點的人，輕視人類所珍視的一切，唯獨鍾愛神聖的生命，因為\BibleQuote{凡有血肉的都似草，人的一切光榮都似草上的花}\BibleRef{伯前 1:24}，他永不會關心這\Quote{今天存在，明天不復存在}的草；研究神聖之道，他不但知道人類生命沒有恆久性，而且整個宇宙也不會永遠平靜運行；他將此世的存在視為異鄉和短暫之物；因為救主說：\BibleQuote{天地要過去}\BibleRef{瑪 24:35}，萬物都必等待改造。只要他\BibleQuote{在這帳棚裡}\BibleRef{格後 5:4}，顯出必死性，被此存在壓迫，他哀嘆寄居的延長；正如聖詠詩人在他的天上詩歌中所說的。的確，那些寄居在這些活帳棚中的人生活在黑暗裡；因此那訓道者呻吟於寄居的持續，說：\Quote{我哀嘆我的寄居延長}，並將沮喪的原因歸於\Quote{黑暗}；因為我們得知黑暗在希伯來語中稱為\OriginalTerm{黑暗}{kedar}。這確實如夜的黑暗籠罩人類，阻止他們看穿這欺騙，認識到活人所最珍視的一切，以及相反的一切，只存在於不反思者的觀念中，本質上空無一物；沒有任何真實事物如出身低微或高貴、榮耀、光輝、古老名望、當下高位、統治他人或臣服。財富與舒適、貧窮與痛苦，以及生命中的其他不平等，在愚昧人看來，以快樂為標準，彼此大不相同。但在更高的悟性看來，它們都一樣；一個不比另一個更有價值；因為生命以同樣的速度通過所有這些對立面奔向終點，在兩類命運中，仍然存在同樣的能力選擇好壞生活，\BibleQuote{藉著左右的武器，藉著惡名與美名}\BibleRef{格後 6:7}。因此，能看清現實的心靈，將沿著自己的道路前行而不轉彎，完成從出生到離世所定的時間；它不被快樂軟化，也不被艱苦擊倒；而是像旅行者一樣，始終前進，對眼前景色不大留意。旅行者的作風是趕往旅程終點，無論是穿過草地和耕種的農田，還是經過更荒野崎嶇的地方；宜人的景色不挽留他們，陰沉之地也不阻攔他們的腳步。同樣，那崇高心靈將直奔自設的目標，不轉向看路上任何事物。它穿越生命，目光卻凝視天上；它是好舵手，把船駛向那裡的某個地標。但較粗俗的心靈向下看；它將精力投注於身體的享樂，正如羊低頭吃草；它為暴食和更低級的快樂而活；它與天主的生命隔絕，與盟約恩許無關；它不承認善，除非是身體的滿足。正是這樣的心靈\Quote{行在黑暗中}，發明了我們生活中一切的邪惡：貪婪、無節制的私慾、無度的奢華、權力慾、虛榮，以及所有入侵人家園的道德疾病之群。在這些惡習中，一個不知如何緊緊連著另一個；其中一個進入，其餘似乎都跟著進來，按自然順序互相拖拽，就像鎖鏈一樣，只要拉動第一個環節，其他環節就不能安靜，甚至最遠端的環節也感受到第一環的運動，藉著毗鄰的關係通過中間環節傳遞過去；人的惡習因著本性如此牢固地相連；當其中一個控制了靈魂，其餘的便隨之而來。如果你想生動描繪這該詛咒的鎖鏈，假定一個人因某種特別的快樂而成為名譽渴求的犧牲品；接著，增加財富的慾望緊隨著名譽渴求而來；他變得貪婪；但只是因為第一個惡習引導他至此。然後這對金錢和優越的貪婪引發對親屬的忿怒、對下屬的驕傲、或對上級的嫉妒；隨後虛偽跟著嫉妒而來；之後是酸腐的脾氣；之後是厭世心；在它們之後是定罪狀態，終結於地獄的暗火。你看這鎖鏈；一切源於一個被珍視的私慾。既然這不可分離的道德疾病之群已經一勞永逸地進入世界，在默感著作的勸誡中，為我們指明了一條唯一的逃脫之路：那就是與包含這連鎖苦難的生活分離。如果我們常駐索多瑪，就不能逃離硫磺火雨；若逃出的人回望它的荒涼，就不能不成為原地生根的鹽柱。我們不能脫離埃及的奴役，除非離開埃及，即這水下生活的世界，渡過那並非紅海，而是這黑色陰鬱的生命之海。但假設我們仍留在這邪惡的奴役中，用師尊的話說，\Quote{真理還沒有使你們自由}，一個追求虛謊並在這世界迷宮中徘徊的人，怎能來到真理面前？一個將自己的存在交給本性鎖鏈的人，怎能逃脫這俘虜？一個例子將使我們的意思更清楚。冬天的洪流，本身湍急，又因水勢上漲，帶走樹木、巨石及沿途一切，對於沿著它的河道生活的人是死亡和危險；但對於遠遠避開它的人，它徒然咆哮。同樣，只有生活在生活動盪中的人才感受到它的力量；只有他才承受那些苦難——本性的河流以麻煩的洪水下降，為了忠於其本性，必將這些苦難帶給在河岸上行走的人。但若人離開這洪流和這些\Quote{狂傲的水}，他將避免成為這生命的\Quote{齒間獵物}，如聖詠接下來所說，並如\Quote{如鳥脫離羅網}，藉著美德的翅膀。因此，這洪流的比喻是成立的；人類生活\Emph{是}一條翻騰喧囂的溪流，向下奔流以達到其自然水平；其中追求的沒有一樣存到尋求者滿足；這溪流帶給他們的一切，靠近、觸及，然後流過；以致在這種湍急流動中的當下瞬間逃脫了享受，因為隨後的急流從他們的視線中將其奪走。因此，我們的利益在於遠離這樣的溪流，免得沉溺於暫時事物而忽略永恆。人如何能永遠保持此處的任何事物，無論他對它的愛多麼熱烈？生命中最珍貴的事物有什麼是永存的？哪一朵青春之花？哪一項力量與美麗的恩賜？財富、名譽或權力？它們都有短暫的花期，然後消散為對立面。誰能持續處於生命的壯年？誰的力量永恆？自然豈不使美麗的花朵比春天的景象更短暫？因為它們按季節綻放，枯萎一段時間後又復甦；再次落葉後又發芽，並將今日的美麗保留到晚期開花。但自然僅在早春的青春期展現人類的花朵；然後就扼殺它；它在年歲的霜寒中消散。其他一切快樂也暫時欺騙肉體的眼睛，然後消失在遺忘的面紗之後。自然不可避免的變化有許多；它們使熱愛者痛苦。只有一條逃脫之道開著：就是不依附任何這些事物，並盡可能遠離這情感和感官世界的交往；或者更確切地說，人走出去，超越自己身體所引發的感覺。這樣，既然他不為肉體而活，就不會受肉體的痛苦所苦。但這就等於只為精神而活，並盡可能模仿精神世界的事工。在那裡，他們不娶也不嫁。他們的工作和卓越在於默觀一切純潔的聖父，並從一切美善的泉源美化自己品格的線條，只要效法祂是可能的。\par

% structure: p0012 source=h2 target=section reason=relative-heading-rank; base=h2

\section{第五章}

现在我们宣告，\OriginalTerm{贞洁}{Virginity}是实现这崇高热望之目标的\Quote{合作者}和助佑。在其他科学中，人们已设计出某些实践方法来培养特定学科；因此，我认为，贞洁正是属神生命科学中的实践方法，赋予人同化于灵性本性的能力。在此历程中，不断的努力是防止灵魂的尊贵因那些感官的爆发而降低——那时心神不再保持属天的思想和向上的凝视，却沉沦于属于血肉的情感。那固着于肉身享乐、忙于纯粹属人渴求的灵魂，如何能以超脱的眼光转向与其同类的理性之光？这种对物质事物邪恶、无知且偏颇的倾向会阻止它。猪的眼睛自然向下，看不见天空的奇观；同样，那被身体向下拖曳的灵魂也再不能仰望上面的美；它不得不凝视那些虽属自然、却是低等而兽性的事物。为能以自由而专注的目光注视天上的欢愉，灵魂将自身从地上转向：它甚至不会分沾世俗生活公认的放纵；它将全部情感能力从物质对象转移到对非物质美的理性默观。身体的贞洁被设计来促进灵魂的这种倾向；其目的是在其中创造对自然情感的完全遗忘；它要防止必要下降到肉身需求之召唤。一旦这样解脱，灵魂就不会冒险——因逐渐习惯沉溺于那些在某种程度上似乎为自然律所容许的事物——而变得对神圣且无玷的欢愉漠不关心和无知。心地纯洁——我们生命的主宰——唯独它能捕获它们。\par

% structure: p0014 source=h2 target=section reason=relative-heading-rank; base=h2

\section{第六章}

我相信，这造成了厄里亚先知的伟大，以及后来在厄里亚的精神与能力中显现的那一位——关于他，有言：\BibleQuote{在妇人所生者中，没有比若翰更大的} \BibleRef{玛 12:11}。如果他们的历史传达了任何其他奥秘教训，那么毫无疑问，这一点藉由他们特殊的生活方式得到了教导：即那心思专注于不可见事物的人必然与一切日常生活事件分离；他对真正之善的判断不会因感官产生的欺骗而混淆和误导。两人自幼便自我流放于人类社会之外，甚至在某种意义上自我流放于人性之外，因为他们忽视了寻常的饮食种类，并寄居在旷野中。各自的需要都由途中遇到的滋养品满足，以便他们的味觉保持简单纯朴，正如耳朵摆脱任何分心的噪音，眼睛摆脱任何游移的目光。如此，他们达到了灵魂的万里晴空般的宁静，并被提升至圣经为每人记载的那种神圣恩宠的高度。例如，厄里亚成了天主在世恩赐的分配者；他有权随意关闭天空的功用以惩罚罪人，或为忏悔者开启它们。若翰确实没有行过任何奇迹；但在他内的恩赐被那洞察人心秘密的那一位宣布为大于任何先知。我们可以推测，之所以如此，是因为两人自始至终如此全心奉献于主，以至于未被任何属世情欲玷污；因为妻儿之爱或任何其他属人的召唤没有干扰他们，他们甚至认为日常的供养不值得焦虑；因为他们显明自己超越任何华服，以所遇之物为满足，一人披山羊皮，另一人穿骆驼毛。我相信，如果婚姻使他们软化，他们就不能达到这种精神的崇高。这不只是简单的历史；它是\Quote{为劝诫我们而写} \BibleRef{格前 10:11}，好使我们能按他们的生活来指引自己的生活。那么，我们借此学到什么？这就是：那渴望与天主结合的人必须像那些圣人一样，使他的心思远离一切世俗事务。被分散在许多渠道中的心思不可能赢得通往认识并热爱天主的道路。\par

% structure: p0016 source=h2 target=section reason=relative-heading-rank; base=h2

\section{第七章}

一个比喻能使我们关于这个主题的教导更加清晰。想象从泉源流出一条溪流，并分岔成许多偶然的水道。只要它这样流淌，就对农业的任何目的毫无用处，因为水的分散使得每条支流变得细小而微弱，因而缓慢。但若有人将这些散漫而广泛散布的小溪重新汇集到一条单一的水道中，他就会有一条充沛而汇聚的溪流，以供应生活所需。人的心智（据我看来）正是如此：只要它的水流向四面八方，散漫于感官的快乐，它就没有任何值得称道的力量趋向\OriginalTerm{真正的善}{Real Good}；但一旦唤回它并将其聚集于自身，以便它开始不再散漫彷徨，而走向与它\OriginalTerm{生来固有的本性活动}{congenital and natural activity}，它就会发现上升到更高事物并把握实在毫无障碍。我们经常看到，水管中的水因约束力而向上喷涌，这种力不许它泄漏；尽管水自然向下流动，同样，人的心智被封闭在\OriginalTerm{习惯性贞洁}{habitual continence}的紧凑水道中，没有任何侧漏，就会因它本来的运动能力而被提升到一种崇高的爱。事实上，它的造物主规定它必须永远运动，停止对它来说是不可能的；因此，当它被阻止在这琐碎之事上使用这种力量时，它必然要冲向真理，因为一切不正当的出口都被封闭了。在许多转弯处，我们看到旅行者能够保持在直路上，当他们知道其他路是错误的，从而避开它们；他们越避开这些错误的方向，就越能保持直路；同样地，心智在远离虚妄时会认识真理。因此，我们刚才提到的那些伟大的先知似乎要教导我们这一点，即不要让自己纠缠于今世努力追求的任何一个对象；婚姻就是其中之一，或者更确切地说，它是一切追逐虚妄之事的\OriginalTerm{首要根源}{primal root}。\par

% structure: p0018 source=h2 target=section reason=relative-heading-rank; base=h2

\section{第八章}

但愿无人以为，我们在此贬低婚姻作为制度。我们深知婚姻并非与天主的祝福无关。但由于人类共有的本能已能为婚姻充分辩护——这些本能出于自然的倾向，促使人选择婚姻的高路以生育子女——而贞洁在某种程度上阻挠了这一自然冲动，因此，正式撰写一篇婚姻劝言实属多余。我们反而以其愉悦作为最有力的卫士为它辩护。然而，尽管如此，或许仍\Emph{有}必要撰写这样的论著，因为有些人歪曲了教会的教导。这些人\BibleRef{弟前 4:2}\Quote{良心已被烙铁烙过}，正如宗徒所说的；这说得非常真切，因为他们离弃圣神的引导，转向\Quote{魔鬼的道理}，心中烙下了溃疡和疱疹，憎恨天主的受造物，称它们为\Quote{污秽}、\Quote{诱惑}、\Quote{有害}等等。\Quote{关于外人，我有什么资格判断呢？\BibleRef{格前 5:12}}宗徒问道。的确，这些人是站在我们奥秘之言得以宣讲的法庭之外；他们并非安置在天主的屋檐下，而是在恶者的隐修院里。他们\Quote{被魔鬼掳去，随从它的旨意\BibleRef{弟后 2:26}}。因此，他们不明白一切德行在于中道，而偏往任何一方都会成为恶习。事实上，把握做得太少与做得太多之间的中点，便区分了恶习与德行。实例将使这更清楚。怯懦和鲁莽是两个公认的相互对立的恶习；一个是自信的不足，另一个则是过度；勇气居于其间。同样，虔诚既不是无神论也不是迷信；否认天主和相信多神同样是不敬。还需要更多例子来阐明这一原则吗？一个人若同时避免吝啬和挥霍，便能通过避开极端而形成慷慨的道德习惯；因为慷慨既不是倾向于随意花费巨大而无用的钱财，也不是在必要开支上精打细算。我们无需逐一详述所有美好的品质。在这一切中，理性已确立德行是两种极端之间的中间状态。因此，节制本身也是一种中间状态，并显然涉及向两边恶习的偏离；也就是说，那些缺乏灵魂的坚韧，在享乐的战斗中轻易被击败，甚至从未踏上德性和节制生活的道路的人，滑向可耻的放纵；而那些越过节制安全之地，超越这一德行的中庸的人，则如同从悬崖上坠入\Quote{魔鬼的道理}中，\Quote{良心已被烙铁烙过}。他们宣称婚姻为可憎恶的，便以此类斥责烙印自己；因为\Quote{如果树是坏的}（如福音所说），\Quote{它的果子也必是坏的}；若人是婚姻之树的枝条和果实，那么对婚姻的斥责便落在那斥责的人身上。因此，这些人如同已被定罪的囚犯；他们的良心被这种非自然教导的鞭痕所覆盖。然而，我们对婚姻的看法是：尽管追求天上的事应是人的首要关切，但若人能以节制和中道利用婚姻的益处，他不必轻视这种服务社会的方式。在圣祖依撒格身上可找到一个例子。他在青春已过、盛年将尽之时娶了黎贝佳，因此他的婚姻并非出于情欲，而是为天主应许给他后裔的祝福。他与她同居直到她生下双生子，然后封闭了感官的通道，完全为那不可见者而活；因为在依撒格历史记载中提到圣祖的\Emph{眼睛昏花}似乎正有此意。但让那些善于解读这些含义的人去那样认为吧，我们继续我们的论述。那么，我们那时在说什么？即：在那些既能忠于神圣的爱又能拥抱婚姻的情况下，没有理由抛弃这一自然安排，并将那尊贵的视为可憎恶的。让我们再次以水和泉源为例。每当农人为灌溉某处而引水，但水到达之前需要一个小出口时，他只允许流出足以供应需求的水量，且这水能轻易与主流重新汇合。若他作为一个缺乏经验、粗心的管水人，把渠道开得太宽，就有整个水流偏离其正道而侧流旁出的危险。同样，如果（因为生命确实需要代代相继）人如此处理这一需要，以致将属灵之事置于首位，并因时日短少而节制性欲、加以约束，那么此人将实现谨慎农人的品格，宗徒也劝勉我们如此。关于偿还这些微不足道的自然债务的细节，他也不会过于计较，但他的长时间祈祷将确保纯洁，这是他生活的基调。他总是害怕因这种放纵而成为血肉之躯；因为天主的圣神不住在血肉之中。性格如此软弱、无法与自然冲动勇敢抗衡的人，最好远离这些诱惑，而不是降入一场超出其力量的战斗。对于他来说，存在着不小的危险：在评估快感时受骗，以至认为除与某种感官情感一同享受的善之外别无他善，并完全远离非物质之爱的享受，从而完全成为属血肉者，只在那里寻求快乐，以致他的品性成为\Quote{爱快乐者}而非\Quote{爱天主者}。并非每人都有此天赋，由于本性的软弱，无法在这些事上把握恰当比例；存在着被远远超越的危险，如圣咏作者所说\Quote{陷入深深的淤泥中}。因此，正如我们的论述所暗示的，我们最好在一生中不经过这些诱惑的考验，以免在合法放纵的借口下，情欲得以进入灵魂的堡垒。\par

% structure: p0020 source=h2 target=section reason=relative-heading-rank; base=h2

\section{第九章}

习俗的确在每件事上都难以抵抗。它具有巨大的吸引和诱惑灵魂的力量。在一个人已经形成固定情感状态的情况下，这种习惯在他内产生了一种对善的想象；没有什么东西天生是卑劣的，但一旦成为风尚，就可能被认为是可欲且可称赞的。试看现今生活在地上的人类。有许多民族，他们的追求并不相同。美丽和荣誉的标准在每一个民族中都不同，每个民族的习俗规范着他们的热情和目标。这种差异不仅出现在那些彼此不重视对方追求的民族之间，甚至在同一民族、同一城市、同一家族中也可以看到；我们也能在这些集合体中看到因习俗感受而产生的许多差异。因此，同一次产痛所生的兄弟，在人生目标上彼此相距甚远。这一点也不奇怪，因为每个人通常不会对同一事物保持同一见解，而是随风尚的影响而改变。不远离我们当前的主题，我们认识一些人，在青春期的早年岁月中表现出对贞洁的热爱；但在他们视为合理和允许的享乐中，他们却以此作为不洁生活的起点，一旦他们接受了这些诱惑，他们一切情感的力量就转向那个方向，并且，再次借用我们关于溪流的比喻，他们让溪流从神圣的渠道冲向低级的物质渠道，在自己内为情欲开辟一条宽阔的道路；以致他们爱情的溪流使更高之路的遗弃渠道干涸，而泛滥于放纵之中。因此，我们认为，软弱的弟兄们最好逃向童贞，如同逃向一座坚不可摧的堡垒，而不是进入人生后果的历程，任凭诱惑在他们身上施行最坏的作为，使自己纠缠于那些通过肉体私欲反对我们心神法律的事物；他们最好考虑，在此他们冒险的不是广阔的田地、财富或其他今生的奖赏，而是那引导他们的希望。一个转向世界、感受其焦虑、并全心致力于取悦人的人，不可能实现师傅的第一条也是最大的诫命：\Quote{\BibleQuote{你应全心、全灵、全力爱你的天主。}\BibleRef{玛 22:37}} 当他将心分裂于天主与世界之间，并将唯独欠祂的爱消耗在人间情感中时，他怎能实现它呢？\Quote{未婚者挂虑主的事，已婚者挂虑世俗的事。} 如果与快乐的战斗似乎令人厌倦，但无论如何，愿所有人鼓起勇气。习惯不会不产生愉悦感，即使是在那些看似最烦躁的人中，通过他们坚持的努力本身而产生；那愉悦将是最崇高、最纯洁的一种；明智的人完全可以喜爱它，而不是让自己因目标的低下而视野狭窄，与那超越一切思想的真正伟大疏远。\par

% structure: p0022 source=h2 target=section reason=relative-heading-rank; base=h2

\section{第十章}

果然，有什么言语能表达那离开了真实美善的拥有而失落的伟大呢？又需要怎样完满的思想力量才能做到！谁能用言语说出心灵所不能把握的事，甚至连轮廓都勾画不出呢？一方面，若一个人保持他心中的眼睛如此清澈，以至于能在某种程度上看见吾\Person{主}真福八端中所许的福乐实现，他将谴责一切人类的言语无力表达他所领悟的。另一方面，若一个人因物质的放纵，使情欲的薄雾像一层薄膜覆盖在他灵魂的锐利视觉上，那么一切表达的力量对他都白费了；因为对于那些毫无感知能力的人，无论你轻描淡写还是夸大奇迹，都是一样的。正如阳光之于一个从出生起从未见过它的人，一切用言语把它传达给他的努力都完全白费；你无法使那光线的光辉透过他的耳朵照射进来；同样，要看见那真正且理智之光的美丽，每个人都需要自己的眼睛；那位因天主灵感的恩赐而能看见的人，则将他的狂喜深藏于意识深处，未表达出来；而那没有看见的人，甚至无法知道自己损失有多大。他怎能知道呢？这美善逃过他的感知，且无法向他呈现；它是不可言说的，不能被描绘。我们尚未学会表达这美丽的独特语言。我们要说的例子在这世界上不存在；为它找一个比较，至少是非常困难的。谁会将太阳比作小火星？或将广阔深渊比作一滴水？那微小的一滴和那一小粒火星，与那深渊和太阳的关系，就如同人所欣赏的任何美丽之物，与那至善面容上的真实之美相比一样——这真实之美是我们超越其他一切美善所瞥见的。能发明什么言语，向那遭受这损失的人显示这损失的宏大呢？在我看来，伟大的\Person{达味}似乎很好地表达了这样做的不可能性。他被圣神的德能提升而超出自我，并在一种有福的狂喜状态中，看见了那无限而不可领悟的美；他看见的程度，正如一个脱去了肉体外壳、单凭思想之力进入精神与理智世界默观的人所能看见的那样；在他渴望说出配得上那景象的话语时，他高呼出那句人人应和的话：\BibleQuote{人人都是说谎者！}我认为这意味着，任何把呈现不可言喻之光的任务托付给言语的人，他实在是真正说谎者；不是因为他憎恨真理，而是因为他表达所思想之物的工具是软弱的。在我们的今生中所遇见的可见之美，无论它以自身的一瞥呈现在无生命的物体或有生命的有机体中，带有某种精选的色彩，都能被我们的审美能力充分欣赏。它可以通过描述被说明并使他者知晓；它可以在语言中被看见，如同在一幅画中被描绘出来。即使这类之美的完美典型也不会难倒我们的概念。但是，当语言找不到其描写的媒介，没有颜色，没有轮廓，没有宏伟的尺寸，没有无瑕的相貌，也没有任何其他的艺术常格时，语言又如何能说明呢？那不可见、无形状、缺乏性质、远离我们用眼睛在物体中所认识的一切的美，绝不可能通过那些只需感官知觉就能把握的特征而被人认识。并不是说我们要对赢得我们爱的这个对象绝望，尽管它似乎高得超出我们的理解。理性越显示我们所寻求之物的伟大，我们就必须越高地提升我们的思想，用那对象的伟大来激发它们；并且我们必须害怕失去我们在那超越的善中的份。事实上，有相当大的危险，由于我们无法将对该物的领悟建立在任何可知的性质上，我们可能会因为它本身的高超和奥秘而完全滑离它。因此，基于思想官能的软弱，我们认为有必要逐步通过感官的认知将它引向那不可见者。我们对这种情况的理念如下。\par

% structure: p0024 source=h2 target=section reason=relative-heading-rank; base=h2

\section{第十一章}

凡以肤浅而不加反思的方式看待事物的人，只观察所遇之物的外表，\Emph{例如}一个人的外表，便不再为其劳心。他们对其身体体积的看法，足以使他们自以为完全认识了他。但具有洞察力和科学精神的心灵，不会将衡量实在的任务仅仅托付给眼睛；它不会停留在现象上，也不将未见之物视为非实在。它探究人灵魂的品性。它考察由身体结构所发展出的那些特征，既看整体，也看个别；先通过分析看个别，然后看那构成主体位格的活生生的整体。关于对美的本性的探讨，我们再次看到，心智未成熟的人，当观察到被某种表面美的光辉所浸染的对象时，便认为那对象本质上就是美的，无论是什么如此先入为主地使他享受视觉的愉悦。他不会进一步深究。但心灵之眼清澈、能审视这些现象的人，则会忽略那些仅仅是美的形式所作用的质料元素；对他而言，这些不过是攀登向那理智之美景象的阶梯，一切其他美的东西都按照各自分有理智之美的程度而获得存在和名称。但对于大多数终生思维如此迟钝的人，我认为，完成这种心智分析、区分质料载体与内在之美、从而把握美的真实本性，是困难的事；如果有人想知道所有虚假有害的美之观念的确切来源，他会发现无非就是这一点：即灵魂的感受官能缺乏那种精确的训练，使之能够区分真美与假美。因此，人们放弃了寻求真美。有些人滑入纯粹的感官享受。另一些人将欲望投向冰冷的金属钱币。还有一些人将美的想象局限于世俗的荣誉、名声和权力。另有一类人热衷于艺术和科学。最卑劣的人以贪饕作为衡量善的标准。但那些远离一切粗俗思想和一切对表象的激情渴求、探索那单纯、非质料、无形之美的本性的人，在诸多可欲对象中做选择时，绝不会犯那样的错误；他绝不会被这些吸引力所误导，以至于看不到其愉悦的短暂性，无法达到对每一个对象的彻底鄙视。因此，这就是引导我们发现美本身的道路。所有其他吸引人之爱的对象，无论多么时尚，无论被何等珍视和热切拥抱，都必须被我们置于下位，因为它们太低、太短暂，不足以运用我们所拥有的爱的能力；并非这些能力要被封闭在我们内不用不动，而是必须首先从一切低级欲望中洁净，然后我们将其提升到感官永远无法企及的高度。甚至对天空之美、耀眼阳光之美，以及任何美丽现象的赞叹，都将停止。在那里注意到的美，不过是指向那至高之美之爱的手，其荣耀是天穹所宣示的，其奥秘是万物所歌颂的。攀升的灵魂，将已把握的一切弃于身后，认为它们对自身需要而言太狭小，如此便能把握那远超诸天的宏伟观念。但是，那些抱负匍匐于地的人，怎能达到这境界？谁能在没有天堂翅膀、尚未因属天召叫而轻盈高远之时，飞入云霄？很少人会对福音的奥秘如此陌生，以至于不知道使人的灵魂升入诸天的唯一载体，就是让自己与那降下的鸽子相似——先知达味也切望其翅膀。这是圣经中用来指圣神能力的寓意名称；可能是因为在那鸟身上找不到一滴苦胆，或者如细心的观察者所说，它不能忍受任何恶臭。因此，凡远离一切苦涩和一切肉身的恶臭，并在前述翅膀上将自己高举于一切低俗世俗野心之上，甚至超越整个宇宙的人，才能找到那唯一值得爱慕的，并使自己变得如同他所触及并进入的美一样美，而且在真实光明的共融中变得明亮发光。那些研究这课题的人告诉我们，夜间在天空中快速相继掠过、有些人称为流星的光芒，不过是空气在特定气流压力下流入天空上层。他们说，当那些气流在太空（以太）中燃烧时，就沿天空划出火迹。同样，环绕大地的空气被一些气流向上推，转化为太空的纯净光辉；人的心灵也离开这泥泞黑暗的世界，受圣神推动，与真实而至高的纯洁接触后变得纯净明亮；在这样的氛围中，它甚至自身发出光来，充满光辉，以致自身成为光，正如我们的主所应许的：\Quote{义人要发光如同太阳。\BibleRef{玛 13:43}}我们甚至在这里也能看到，例如镜子、水面或任何能反射光的光滑表面；当它们接收阳光时，自己也发光；但如果任何污渍破坏了其纯净光亮的表面，它们就不能这样。如果我们离开这大地的黑暗气氛而居于高处，亲近基督的真光，我们也将成为光；如果那照亮黑暗的真光也降临到我们身上，我们就是光，如主在某处对门徒所说的；除非罪的污秽漫延到我们心上，使我们光明变得暗淡。或许这些例子逐渐引导我们发现，我们能够被改变为比自己更好的存在；并且也证明，灵魂与不朽之天主的这种合一，除了她通过童贞状态达到最大可能的纯洁之外，别无他法——这种状态因相似于天主，使她能把握与她相似的那一位，同时她将自己如镜子般置于天主的纯洁下，并在触碰和观看一切美的原型时塑造自己的美。取一个足够坚强的人格，能远离一切属人的事物——人物、财富、艺术和科学的追求，甚至远离在道德实践和立法中被视为正确的事物（因为在所有这些中，对美的理解仍有错误，感官是标准）；这样的人只会成为那美的热恋者，这美本身是唯一的根源，不是在某一特定时间或相对而言，而是从自身、通过自身、在自身内是美的，不是一时是美、下一刻就不美，超越一切增加和添加，不能改变和变化。我敢肯定，对于那已清净其存在的一切能力、摆脱各种恶习的人，本质之美——一切美和一切善的源泉——将成为可见的。视觉之眼，清除了使人目盲的体液，能清晰辨认甚至遥远天空中的物体；同样，灵魂因其纯洁而获得接受那光的能力；真正的童贞、真正的贞洁热忱，其最终目标无非就是：藉此能力看见天主。事实上，没有人如此心智盲目，以至于无需告知就不明白\Emph{那}一点：即宇宙的天主是唯一绝对、最初且无与伦比的美和善。或许所有人都知道那点；但如所料，有些人还想尽可能发现一个过程，使我们可以实际被引导到它。神圣的经卷充满了这样的指导；此外，许多圣人像明灯一样，以其生命的光辉照耀那些\Quote{与天主同行}之人的道路。但每个人可以从默感著作的两约中为自己收集丰富的达此目标之提示；先知和法律充满了它们；福音和宗徒的传统也是如此。我们自己在追随那些默感之言的思想时所推测的，便是如此。\par

% structure: p0026 source=h2 target=section reason=relative-heading-rank; base=h2

\section{第十二章}

这个有理性和智力的受造物，即人，既是神圣不朽之理智的化工，又是祂的肖像（因为在创世记中论到他这样写着：\BibleQuote{天主照自己的形象造了人}\BibleRef{创 1:27}），我说，这个受造物在其最初被造的过程中，并没有将情欲和死亡的可能性与他自己本性的本质结合在一起。事实上，如果肖像所反映的美与原型之美有丝毫相悖，那么关于肖像的真理就绝不可能得到维护。情欲是在后来，即在人的最初构造之后才被引入的，其方式如下。如所说，人是统治万有之能的肖像和模样，因此在自由意志方面，他也保持了与那意志统御万有的那一位的相似。他不受任何外在必然性的奴役；他对所悦纳之物的情感只取决于他自己的私人判断；他可以自由选择他所喜欢的任何事物；因此他是一个自由的行动者，尽管他被诡计所蒙蔽，从而招致了如今压在人身上的那场灾难。他本人成了邪恶的发现者，但他所发现的并非天主所造的；因为天主并未造死亡。事实上，在某种程度上，人自己成了邪恶的制造者和工匠。所有拥有视觉能力的人都能同样享受阳光；任何人若愿意，都可以闭上眼睛拒绝这种享受：在这种情况下，并非太阳隐退而产生黑暗，而是人自己在眼睛和阳光之间设置了一道屏障；即使在闭眼时，视觉能力也不能保持不活动，因此这种活动的视觉必然变成一种活动的黑暗，这黑暗是因他停止观看的自由意志行为而在人内升起的。又比如，一个人在为自己建造房屋时，可能忽略在其中留出任何光线的入口；他必然会处于黑暗之中，尽管他是自愿切断自己与光的联系的。同样，地上的第一个人，或者更确切地说，那在人内生出邪恶的那一位，原本可以选择那在万物本性中环绕着他的善与美；然而他却故意违背这本性为自己开辟了一条新的道路，并在远离德行的行为（这是他自己的自由行为）中创造了邪恶的习惯。因为，请注意，世界上并没有一种与意志无关并能独立于意志而在实体中发现的邪恶。天主的每一受造物都是好的，没有一样是\BibleQuote{可弃绝的}；天主所造的一切都\BibleQuote{很好}。但是，正如我们所描述的，犯罪的习惯进入了人的生命，并带着致命的迅速，从那个小小的开端蔓延成如今这无穷的邪恶。然后，那原本模仿原型之美的灵魂的神圣之美，就像优质钢铁被邪恶的铁锈染黑一样，不再保持其惯常本质的荣耀，反而被罪的丑陋所毁容。这如此伟大而珍贵（正如圣经所称的）之物，即人，已经从其骄傲的天赋权利中堕落了。正如那些滑倒并重重摔入泥沼的人，整个面部都被污泥涂抹，以至于最亲近的朋友也认不出他们来；同样，这个受造物也跌入了罪的泥沼，失去了作为不朽天主肖像的福分；反而为自己穿上了另一种可朽坏且污秽的相似物；而理性劝告他通过在此召叫的洁净水中洗去这层污垢来再次除去它。一旦除去尘世的包裹，灵魂的美便会再次显现。现在，除去一种陌生的附着物，就等于回归那熟悉的、自然的状态；然而，这样的回归只有通过再次成为我们在起初被造时的样子才能实现。事实上，这种与神性的相似根本不是我们的工作；它不是人的任何能力的成就；它是天主在我们出生那一刻就赐予我们本性的伟大恩赐；人的努力只能达到清除罪的污秽，从而使被掩埋的灵魂之美再次闪耀。我想，这个真理在福音中得到了教导，当我们的主对那些能领会智慧在奥秘之下的言语的人说：\BibleQuote{天主的国就在你们中间}\BibleRef{路 17:21}。这句话指出，神性的善并非与我们本性分离的东西，也不是远离那些愿意寻求它的人；事实上，它就在我们每个人内，当它被生活的忧虑和快乐窒息时，确实被忽视和未被注意，但每当我们能够将我们意识思考的能力转向它时，就会重新找到它。如果我们所说的话需要进一步证实，我认为它会在我们的主寻找失落的达玛的比喻中找到。那里的思想是，守寡的灵魂从其他德行（比喻中称为达玛）中得不到任何益处，即使它们全都完好，如果那一个不在这其中。因此，比喻建议先点上蜡烛，无疑象征我们的理性，它照亮隐藏的原则；然后，在自己家中，即在自己内，寻找那枚失落的钱币；比喻无疑暗示那钱币是我们君王的肖像，尚未完全失落，只是藏在污垢之下；而这里我们必须要理解肉体的不洁，通过生活的谨慎被扫除和净化后，使我们寻找的对象清晰可见。然后，比喻的意思是，找到这肖像的灵魂自身为之欢喜，并且她的邻人也被请来与她同享这喜乐。诚然，所有与灵魂同住的力量（比喻为此时称她们为邻人），当那大能君王的肖像最终在其全部光辉中显现时（即每一位人心的塑造者印在这枚达玛上的肖像），将转而沉浸在那神圣的喜乐和欢庆中，并凝视那被寻回者的不可言喻的美。她说：\BibleQuote{你们与我一同欢乐罢，因为我找到了我所失落的达玛。}邻人，即灵魂的惯常力量，包括理智的和欲求的，悲伤和愤怒的情感，以及在灵魂中被辨识的所有其他力量，在那庆祝寻得天上达玛的欢乐节庆中，也被恰当地称为她的朋友；当她们都注视那美与善，并为天主的荣耀而行一切事，不再作罪的工具时，她们理当在主内一同欢乐。\BibleRef{罗 6:13}倘若这就是失物寻得的教导，即我们应该从肉身包裹的污秽中恢复天主的肖像，使其回到原始状态，那么就让我们成为第一个人最初呼吸时的那样子。那是什么样子呢？那时他还没有穿着死皮的外衣，却能毫不畏缩地凝视天主的容颜。他尚未凭味觉或视觉来判断什么是可爱的；他只在主那里找到一切甘甜；他使用所赐予的配偶只是为了这种喜乐，正如圣经所说的，直到他被赶出乐园，并且她因被诱骗所犯的罪被判处生育之苦楚时，\BibleQuote{他还不认识她}\BibleRef{创 4:1}。我们这些在第一位祖先中被逐出的人，被允许通过我们失去乐园的那些相同阶梯返回到我们最初有福的状态。这些阶梯是什么？诡计多端的快乐开始了堕落，快乐之后是羞耻和恐惧，甚至不敢继续留在他们的创造者眼前，以致他们藏身于树叶和树荫中；之后他们用死兽的皮遮盖自己；然后被派遣到这个瘟疫横行、苛刻的大地上，在这里，作为对必死的补偿，婚姻被建立。\BibleRef{创 3:16}如果我们注定要\BibleQuote{离去，与基督同在}\BibleRef{斐 1:23}，就必须从离别之路的终点（离我们最近的起点）开始；正如那些远离家乡的朋友的人，当他们转身要回到出发地时，首先离开的是他们在远行结束时到达的地区。婚姻，是我们与乐园生活分离的最后阶段；因此，正如我们的论述所提示的，婚姻是首先要离开的；它是我们为归向基督而离别的第一站。其次，我们必须从一切焦虑的土地劳作中退隐，那是人在犯罪后所受的束缚。接着，我们必须脱去那些遮盖我们赤裸的外衣，即死皮的衣袍，也就是肉身的智慧；我们必须弃绝一切暗中所行的可耻之事，\BibleRef{格后 4:2}不再用这苦涩世界的无花果叶遮盖自己；然后，当我们撕掉这易朽生命叶子的涂层后，必须再次站在我们的创造者眼前；并排斥一切味觉和视觉的幻觉，只以天主的诫命为我们的向导，而不是那吐毒液的蛇。那条诫命是：只触摸那善的，不品尝那恶的；因为不能忍耐停留在对恶的无知中，只会成为实际邪恶漫长系列的开端。为此，我们的原祖被禁止同时把握对善的对立面的知识和对善本身的知识；他们应当避免\BibleQuote{认识善恶}\BibleRef{创 2:17}，而只享受纯粹的善，不掺杂一丝邪恶；而享受\Emph{那}，在我看来，无非是永远与天主同在，并持续不断地感受这种喜乐，不被任何能使我们与它分离的事物所掺杂。甚至有人可以大胆地说，这可能是人再次被提入乐园的道路，离开这个处于邪恶中的世界，进入那乐园，\Person{保禄}在那里看见了人不可言说、人不可谈论之事。\BibleRef{格后 12:4}\par

% structure: p0028 source=h2 target=section reason=relative-heading-rank; base=h2

\section{第十三章}

然而，既然乐园是活灵的家园，不容纳那些死于罪恶的人，而我们另一方面是属血肉的，屈服于死亡，被卖给了罪恶，那么，一个受死亡帝国管辖的人，怎么可能永远居住在这片全是生命的土地上呢？有什么方法可以摆脱这种辖制呢？在这里，福音的教导也是极其充分的。我们听到主对尼苛德摩说\BibleQuote{从肉体生的就是肉体；从圣神生的就是神。}我们也知道，肉体因罪恶而屈服于死亡，但天主的圣神既是不腐朽的，又是赋予生命的，且是永不死灭的。正如在我们肉体的出生时，一同来到世界的是一种再次化为尘土的可能性，显然，圣神也将赋予生命的可能性赐予祂自己所生的子女。那么，从这些话中得出什么教训呢？这就是：我们应使自己脱离这肉体生命，它有一个不可避免的随从——死亡；并且我们应寻求一种不带来死亡的生活方式。而童贞的生活正是这样的生活。我们将补充一些其他内容，以表明这是多么真实。人人都知道，有死之躯体的繁衍是两性交合的工作；然而对于那些与圣神结合的人，这后一种交合所产生的是生命和永福，而不是儿女；宗徒的话正合适他们的情形，因为这样的孩子的快乐母亲\BibleQuote{必因生育而得救 \BibleRef{弟前 2:15}}；正如圣咏作者在其神圣诗歌中感恩地呼喊：\BibleQuote{祂使不生育的妇女安居家中，成为快乐子女的母亲。}真是一位快乐母亲！这位童贞的母亲，借着圣神的运作，怀有了不死的子女，却因她的谦逊而被先知称为不生育的。那么，这种比死亡更有力量的生命，对于那些思考的人来说，是更可取的。肉体将子女带入世界——我这样说无意冒犯——既是生命的起点，也是死亡的起点；因为从出生那一刻起，死亡的过程就开始了。但是那些借着童贞停止了这个过程的人，在自己内划定了死亡的边界，并以自己的行为阻止了死亡的进展；事实上，他们使自己成为生命与死亡之间的边界，也是一道阻碍死亡的屏障。因此，如果死亡不能越过童贞，反而在那里发现其力量被阻止和粉碎，那就证明了童贞比死亡更为强大；并且那身体被恰当地称为不朽的，它不为这个世界服务，也不肯成为不断死亡之受造物的工具。在这样的身体里，那从第一个人到所过的童贞生活之间，长久不断的腐朽与死亡的进程被打断了。确实，只要人类借着婚姻也在工作，死亡就不能停止其工作；他随着所有先前的世代行走生命之路；他与每个初生的婴儿一同开始，并陪伴其到终点：但在童贞中他发现了一道屏障，要跨越它是不可能的壮举。正如在玛利亚——天主之母的时代，那从亚当直到她的时代实行统治的（死亡），当他来到她面前，将他的力量冲击她童贞的果实如同冲击岩石时，就在她身上被粉碎了；同样，在每个借着童贞的护卫度过此生肉体的灵魂中，死亡的力量在某种程度上被击碎并消除，因为他找不到可以刺入他的毒刺的地方。如果你不向火中投入木头、稻草或干草，或任何它可以燃烧的东西，火本身就没有持续的力量；同样，如果婚姻不向死亡提供材料，并为这个刽子手准备好牺牲品，死亡的力量就无法继续运作。如果你还有任何疑惑，请思考死亡带给人类的那些痛苦的实际名称，这些在本论的前一部分已详述。它们从何获得意义？\Quote{寡居}、\Quote{孤苦}、\Quote{丧子}，如果婚姻没有先行，它们会成为悲伤的主题吗？不，所有婚姻中宝贵的幸福、狂喜和安慰，都以这些痛苦的悲伤告终。剑柄光滑顺手，外表抛光闪亮；似乎与手的轮廓相贴合；但另一部分是钢铁和死亡的器具，看起来可畏，遭遇起来更可畏。婚姻就是这样。它为感官的把握提供了光滑的快乐表面，如同一个罕见的抛光精美工艺的剑柄；但当人拿起它，握在手中时，他发现与之相连的痛苦也在他的手中；它便成了哀悼和损失的工作者。正是婚姻展现了令人心碎的景象：孩子们在年幼时被遗弃孤苦，成为强者的猎物，却往往因不知即将到来的灾祸而对不幸微笑。寡居的原因若非婚姻又是什么？而远离婚姻将带来对这些心灵之沉重税赋的豁免。我们还能期待不同吗？当起初对犯罪者所宣布的判决被废止时，母亲的\BibleQuote{痛苦 \BibleRef{创 3:16}}也不再\BibleQuote{增多}，也没有\BibleQuote{痛苦}预示着人的出生；那么所有灾祸都从生命中移除，\BibleQuote{泪水从每张脸上被擦去 \BibleRef{依 25:8}}；怀孕不再是罪过，生育也不再是罪恶；出生将不再是\BibleQuote{出于血统}，或\BibleQuote{出于人的意愿}，或\BibleQuote{出于肉体的意愿}，而是只出于天主。每当任何人在活泼的心中怀有圣神的全部完整，并产生智慧、正义、圣化和救赎时，这总是在发生。任何人都可能成为这样一个儿子的母亲；正如我们的主说\BibleQuote{那承行我旨意的，就是我的兄弟、姊妹和母亲 \BibleRef{玛 12:50}}。在这样的分娩中，哪里有死亡的位置？实际上，在其中死亡被生命吞灭了。事实上，童贞的生命似乎是未来世界真福的实际表征，它本身显示了那么多为我们保留在那里的所期望的真福临在的迹象。为使这一陈述的真实性得以被感知，我们将如此验证。首先，一个人既已这样一次而永远地死于罪恶，今后就为天主而生活；他不再结出死亡的果实；并且在他能力范围内，已结束了内在的肉体生命，从此等候伟大天主显现的所期望的真福，避免通过中间的后裔拉开自己与天主来临之间的距离。其次，因为他在现世生活中就享受了我们复活的所有真福结果的某种精致的荣耀。因为我们的主宣告，我们复活后的生命将如同天使的生命。天使本性的特点就是与婚姻无涉；因此，这预许的真福已为他所领受，他不仅将自己的荣耀与圣人的光环相融合，而且借着生活的无玷，如此模仿了这些无体之体的纯洁。那么，如果童贞能为我们赢得这样的恩惠，什么样的言语能表达对如此伟大恩宠的赞叹呢？还能找到什么灵魂的恩赐如此伟大而珍贵，以至于不与这完美相比而感到逊色呢？\par

% structure: p0030 source=h2 target=section reason=relative-heading-rank; base=h2

\section{第十四章}

但若我们终于领会这恩宠的圆满，也必须明白其必然结果：即这并非终止于制服身体的单项成就，而是其意向触及并渗透一切属于或被认为是灵魂的正确状态。那在童贞中依附着真正新郎的灵魂，绝不会仅仅远离一切肉体污秽；她将使这种戒除成为她纯洁的开端，并将这种免于失败的安全同样带入她道路上的一切事物。她担心，因着一颗过分偏狭的心，若在任一方向接触邪恶，就会为不忠的最小软弱提供机会（设想如此情况：但我将重新开始我原打算说的话），那依附其主人以致与祂成为一灵，并借着婚姻生活的盟约将她全心、全力的爱只押注于祂的灵魂——这灵魂绝不会再犯任何其他违反救恩的过犯，如同她不会因依恋淫乱而危及与祂的结合；她知道所有罪恶之间有一种共同的不洁亲缘，若她只被一个玷污，就再不能保持无瑕。用一个比喻说明我们的意思。假设池水全部平静无波，没有任何扰动破坏那处的安宁；然后一块石头投入池中；这一处的运动将扩展至整个水面，当石头的重量将其带向池底时，其周围激起的波纹成圈扩散到其他水面，如此经过所有中间的扰动一直推到水边，整个水面被这些波纹搅动，感受到深处的运动。同样，灵魂广阔的宁静与平安被一股入侵的激情所扰，并因一部分的伤害而受影响。那些研究过此问题的人也告诉我们，德行彼此并非分离，若不掌握作为其余德行基础的原则，就不可能掌握任何一项德行的原则；在品格中展现出一种德行的人必然展现出所有德行。因此，相反地，我们道德本性中任何事物的败坏将扩展至整个德行生活；的确，正如宗徒告诉我们，整体受部分影响，\BibleQuote{若一个肢体受苦，所有肢体一同受苦，}\BibleQuote{若一个得荣耀，所有肢体一同喜乐。}\par

% structure: p0032 source=h2 target=section reason=relative-heading-rank; base=h2

\section{第十五章}

但我们生活中转向罪恶的道路不可胜数；圣经以多种方式述说其数目。\BibleQuote{迫害我、加害我的人很多，}\BibleQuote{从高处攻击我的人很多，}以及其他许多类似经文。的确，我们可以断言，像奸夫那样图谋毁坏这真正可敬的婚姻、玷污这无玷之床的人很多；若我们必须一一列出，我们将愤怒、贪婪、嫉妒、报复、敌意、恶意、仇恨以及宗徒归于与纯正道理相悖之列的一切归入这淫乱的精神。现在设想一位贵妇，美丽超群，因此与国王成婚，但因美貌而被放荡的情人围困。只要她对那些诱奸者仍感愤慨并向其合法丈夫控诉，她就保持贞洁，眼中只有她的新郎；放荡者找不到进攻她的立足点。但若她听从了其中一人，她对其余人的贞洁也不能免除惩罚；只要她允许那人玷污了婚床，就足以定她的罪。同样，其生命在于天主的灵魂不会在任何那些以美丽外表欺骗她的事物中找到快乐。若她在一时的软弱中允许污秽进入内心，她就自己破坏了属灵婚姻的盟约；正如圣经告诉我们：\BibleQuote{智慧不会进入恶意灵魂\BibleRef{智 1:4}。}简言之，可以真实地说，良善的丈夫不能与易怒、怀恶意或怀有任何与和谐相悖之情的灵魂同居。尚未发现任何调和本性敌对且毫无共同点之事物的途径。宗徒告诉我们：\BibleQuote{光明与黑暗没有相通\BibleRef{格后 6:14}，}或正义与不法，简言之，我们观察并称为天主本性本质的一切品德，与在恶中观察到的所有相反品德之间没有相通。鉴于相互排斥者之间的任何结合都是不可能的，我们明白邪恶的灵魂与良善相伴是疏远的。那么从这个教训中得出的实际教导是什么？贞洁而深思的童贞女必须断绝任何可能以任何方式给灵魂带来感染的情感；她必须为娶她的丈夫保持纯洁，\BibleQuote{没有瑕疵、皱纹或任何这类事。}\par

% structure: p0034 source=h2 target=section reason=relative-heading-rank; base=h2

\section{第十六章}

只有一条正道。它狭窄而受束缚。它没有向左或向右的转弯。无论我们如何离开它，都会同样有迷失而无可救药的危险。既然如此，许多人养成的习惯必须尽可能改正；我指的是那些人，他们一边奋力对抗低级的快乐，一边却仍在追逐以世俗荣誉和满足权力欲的职位为形式的快乐。他们就像某个渴望自由的仆人，却不努力挣脱奴役，只是换了个主人，以为自由就在于这种改变。但所有奴隶都是奴隶，即使他们不一直受同一主人统治，只要有任何形式的统治和力量凌驾于他们之上，他们都是奴隶。另有一些人在长期对抗所有快乐之后，轻易地在另一个战场屈服，那里有相反的情感介入；在他们极其严谨的生活中，很容易陷入忧郁、恼怒、怀恨以及一切与愉快情感直接对立的东西；他们很不愿意从中解脱。每当任何一种情感而非理性的德行控制了一个人的生活时，总会发生这种情况。因为主的诫命极其明亮，以至能\Quote{照亮眼睛}甚至\Quote{愚钝人}，宣告唯独天主是善。但天主既不是痛苦也不是快乐；既不是怯懦也不是勇敢；既不是恐惧也不是愤怒，也不是任何支配无修养灵魂的情感，而是如同宗徒所说的，祂是真智慧与圣化、真理、喜乐与平安，以及一切类似的事。如果祂是这样，那么一个被完全相反情感所支配的人怎能说是依附于祂呢？一个人认为当他成功地逃脱了某一特定情欲的统治后，就会在其相反的方面找到德行，这难道不是缺乏理性吗？例如，认为当他逃脱了快乐后，就会在让痛苦占据他中找到德行；或者当他努力抵御愤怒后，就会在恐惧中蜷缩。无论我们以这种方式还是那种方式错过德行，或者更确切地说，错过作为德行总和的天主自身，结果都一样。以身体极度虚弱为例：可以说，这种虚弱的悲伤是一样的，无论原因是过度节食还是暴饮暴食；两者未能适时停止都导致相同的结果。因此，守护灵魂生命和健康的人会安守真理的适度；他会继续不触及那两种伴随德行的相反状态。这不是我个人的教导，而是来自天主的圣口。它清楚地包含在吾主对祂的门徒所说的那段话中，即他们如同在狼群中的羊，但不仅是鸽子，而要也在性情中有些蛇的成分；这意味着他们既不应过度实践那看似可称赞的单纯，因为这种习惯近乎疯狂，也不应视大多数人钦佩的机巧为德行，而不以相反的成分来缓和；他们实际上要通过混合这两种看似相反的东西来形成另一种性情，去除一方的愚蠢和另一方的邪恶狡猾；这样，从一个单一美丽的品格中，创造出单纯的目的与精明的结合。祂说\Quote{要如同蛇一样机警，如同鸽子一样纯洁}。\par

% structure: p0036 source=h2 target=section reason=relative-heading-rank; base=h2

\section{第十七章}

吾主当时所说的，应成为每个人生活的普遍准则；尤其应成为那些透过贞洁之门更接近天主之人的准则，即他们决不可在某一方向追求成全的同时，在相反的另一方面暴露未加防备的缺口；而应在所有方面实现善，从而确保他们的圣洁生活在一切事上无懈可击。一名士兵不会只在某些部位武装自己，而让身体其余部分暴露无保护。如果他在那未受保护的地方受了致命伤，那部分盔甲又有何用？同样，谁会称那因意外缺失了构成整体美所必需之要素的容貌为完美？残缺部分的畸形破坏了未受损部分的优雅。福音暗示，那着手建造一座塔楼，却将全部劳力用于地基而从未达到完工的人，是可笑的；我们从塔楼的比喻中还学到什么，除非是努力完成每一项崇高目的的终局，在天主所有诫命的多重结构上完成祂的工程？一块石头毕竟不是整个塔楼的建筑，同样，仅守一条诫命也不能使灵魂的成全达到所需的高度。地基固然必须首先铺好，但宗徒说：\BibleRef{格前 3:12}，在其上必须建造金子和宝石的工程；因为先知所呼喊的也正是如此：\Quote{我喜爱祢的诫命胜过黄金和许多宝石}。让贞洁之爱作为德行生活的基础；但在此之上，应建造每一种德行的成果。如果贞洁被认为是一件极其珍贵且具有神圣外貌的事（确实如此，人也这样认为），然而，如果整个生活不与这完美的音符和谐，反而被灵魂随后出现的不和谐所破坏，那么这事就变成了\Quote{金环戴在猪鼻上}或\Quote{珍珠被猪践踏}。关于这一点，我们已经说得够多了。\par

% structure: p0038 source=h2 target=section reason=relative-heading-rank; base=h2

\section{第十八章}

假设有人认为，在他的生活与其中任何一个情况之间缺乏和谐是无关紧要的，那就让他通过观察一个家庭的管理来学习我们这句格言的意义。一家之主不会容许在屋里看到任何杂乱或不雅，比如床榻翻倒、桌上堆满垃圾、贵重器皿被扔进肮脏的角落，而那些用于卑贱用途的器物却被推到显眼处，使来访的客人看到。他把一切都整齐地安排在最合适的地方，使它们摆放得最得体；然后他就可以迎接客人，不必担心打开屋内接待他们会感到羞愧。我想，同样的责任也落在我们\Quote{帐幕}的主人——理智身上；它必须安排我们内在的一切，并把灵魂的每一种官能——造物主将它们塑造成我们的工具或器皿——用于恰当而高尚的用途。我们现在要详细说明一个人如何管理自己目前处境中的生活以求进步，希望没有人会指责我们琐碎或过于细致。因此，我们建议，将爱的激情置于灵魂最纯洁的圣所，作为我们所有恩赐中选出的初果，完全奉献给天主；一旦这样做，就要保持它不受任何世俗污秽的沾染和玷污。然后，愤慨、愤怒和憎恨必须像看门狗一样，只对攻击性的罪孽才被激起；它们必须只对那潜入劫掠神圣宝库的窃贼和仇敌才遵循自然的冲动，而那贼来只是为了偷窃、撕毁和毁灭。勇气和信心是我们手中的武器，以抵御恶人突发袭击和进攻。望德和忍耐是我们要倚靠的杖，每当我们因世俗的考验而疲倦时。至于悲伤，我们必须储备一些，以备在为自己的罪忏悔时使用；同时相信它除了服务于那目的外，从无用处。正义将是我们正直的准则，保护我们不在言语或行为上失足，并指导我们调配灵魂的官能，以及对所遇的每一个人给予应有的尊重。对利益的贪爱，这在每个灵魂中是一个巨大而无可估量的因素，一旦被用于对天主的渴望，就会祝福拥有它的人；因为他会在该激烈之处激烈。明智和审慎将是我们最好利益的顾问；它们会安排我们的生活，使我们永远不会因任何轻率的愚行而受苦。但如果一个人不将上述灵魂的官能用于其正当用途，反而颠倒了它们本来的目的；假设他把爱浪费在最卑贱的对象上，只对自己的亲属积攒仇恨；假设他欢迎不义，只对父母逞威，只在荒谬之事上大胆，把希望寄托于虚空，将审慎和明智从同伴中驱赶，以贪饕和愚行作为主母，并以同样方式利用其他一切机会，那他确实会成为一个奇怪而非自然的角色，其程度超乎任何人所能描述。如果我们能想象有人把盔甲都穿反了，头盔倒过来遮住脸而羽毛向后翘，把脚伸进胸甲，把胫甲套在胸前，把本该在左边的都换到右边，\Emph{反之亦然}，这样一个\OriginalTerm{重装步兵}{hoplite}在战场上会遭遇什么，那我们就大概能想象那种命运——等待着那个因判断混乱而颠倒灵魂官能正当用途的人。因此我们必须为所有情感提供这种平衡；真正的理智节制自然能够提供它；如果有人要给这种节制下一个准确定义，他可以断然宣称，它等同于我们凭借明智和审慎对灵魂的每一种情感进行有秩序的掌控。此外，灵魂中的这种状态将不再需要任何费力方法来达到崇高而属天的现实；它将以最大的容易实现那先前似乎不可企及的事；它将把握其寻找的对象，作为拒绝相反诱因的自然结果。走出黑暗的人必然在光明中；没有死的人必然活着。的确，一个人若不是白白领受了灵魂，他必定会走在真理的道路上；用以防范错误的审慎和知识本身就是在正路上确实的引导。被释放的奴隶不再服侍旧主，就在他们成为自己的主人的那一刻，就将一切思想转向自己；同样，我想，从服侍身体中解放出来的灵魂立刻意识到自己内在的力量。但这种自由，正如我们从宗徒所学的（\BibleRef{迦 5:1}），在于不再被奴役的轭所束缚，也不再像逃犯或罪犯一样被婚姻的锁链再次捆绑。但我必须回到开头时所说的话：这种自由的完美并不只在于不结婚这一点。没有人应该以为贞洁的赏报如此微不足道且轻易可得，好像一件小小的肉体行为就能解决如此重要的事情。但我们已经看到，凡犯罪的，就是罪恶的奴隶（\BibleRef{若 8:34}）；因此，在任何行动或任何习惯中向邪恶的堕落都使人成为奴隶，而且是烙了印的奴隶，用罪恶的鞭笞使他满身伤痕和灼痕。因此，那抓住所有贞洁的超凡目标的人，理应在各个方面忠于自己，并在生活的一切关系中同样彰显他的纯洁。如果需要任何受默感的话语来辅助我们的恳求，真理本身足以确证真理，当它以福音比喻的隐晦含义教导这种教训时：渔夫凭技巧把好的可吃的鱼和坏的有毒的鱼分开，这样对好鱼的享用就不会因任何坏鱼混入\Quote{器皿}中而被破坏。真正节制的工作也是如此；从一切追求和习惯中，选择纯洁和有益的东西，拒绝任何看来可能无用的，并让它回到普遍世俗的生活中，在比喻的意象中被称为\Quote{海（\BibleRef{玛 13:47}），}。圣咏作者在阐述完全告解的教义时，也将这种不安受苦的动荡生活称为\Quote{水淹至灵魂，}、\Quote{深渊之水，}和\Quote{暴风}；在这海中，一切叛逆的思绪都像埃及人那样，如石头般沉入深渊（\BibleRef{出 15:10}）。但我们内一切为天主所珍爱、并对真理有锐利洞察力的（在叙述中称为\Quote{以色列}），却唯独能像经过干地那样渡过那海，从不被生命波涛的苦咸触及。因此，作为预像，在法律领导之下（因为梅瑟是那将要来临的法律的预像），以色列干爽地渡过那海，而追随其后的埃及人却被淹没。各人按自己所带的性情而行：一人轻松行走，另一人被拖入深水。因为美德是轻盈飘浮的东西，所有按她方式生活的人都\Quote{如云飞翔，}正如依撒意亚所说，\Quote{又如带着幼鸽的鸽子}；而罪却是沉重之事，\Quote{坐，}另一位先知说，\Quote{在一塔冷通铅上。}然而，如果有人觉得对这段历史的这种解释牵强而不适用，并且认为红海的奇迹不是为我们的益处而写的，那就让他听宗徒的话：\Quote{现在这一切事发生在他们身上，作为预像，并且被记载下来，为劝诫我们。}\par

% structure: p0040 source=h2 target=section reason=relative-heading-rank; base=h2

\section{第十九章}

但是，除了其他事情，女先知米黎盎的行为也引起了我们的这些推测。刚一渡过海，她手中就拿起一面干燥而响亮的手鼓，带领妇女们跳舞。\BibleRef{出 15:20} 手鼓可能暗示童贞，因为米黎盎最先成就了童贞；我确实相信她是天主之母玛利亚的预像。正如手鼓因完全不含水分、达到最大程度的干燥而发出响亮的声音，同样，童贞在人间有清晰响亮的声誉，因为它排斥了纯肉体生命的活力汁液。因此，米黎盎的手鼓是死物，而童贞是肉情欲的致死，那么米黎盎是贞女的可能性或许离概率的界限并不很远。然而，我们只能猜测和推测，无法清楚证明事实如此，即女先知米黎盎带领了一群童贞女舞蹈，尽管许多学者明确断言她未婚，因为历史既未提及她的婚姻，也未提及她做母亲；而且，如果她有丈夫，她一定不会被称为\Quote{亚郎的姊妹}（\BibleRef{出 15:20}），而会以丈夫的名字被称呼和认识；因为女人的头不是兄弟，而是丈夫。但是，如果在这样一个民族中，做母亲被追求、被列为祝福、被视为公共责任，而童贞的恩宠却被视为珍贵，那么我们这些不\Quote{判断}天主的祝福\Quote{凭肉身}的人，应该对它持有怎样的情感呢？事实上，在天主的圣谕中已经启示了，在何种情况下怀孕和生育是好事，天主圣人们渴望的是何种生育能力；因为依撒意亚先知和神圣的宗徒都已清楚明确地阐明了这一点。前者呼喊说：\Quote{主啊，因了对祢的敬畏，我怀了孕；}后者夸耀自己是所有大家庭的家长，生养了整座城市和整个国家；不仅是格林多人、迦拉达人——他通过自己的劳苦为主塑造了他们，而且是从耶路撒冷直到依里黎苛的广阔区域；他的儿女充满世界，是他在基督内借着福音所\Quote{生}的。同样，圣童贞的子宫，曾服务于无玷的生产，在福音中被称颂为有福的；因为那生产并未取消童贞，童贞也未妨碍如此伟大的生产。当\Quote{救恩之神}——如依撒意亚所称——被生育时，肉性的意愿是无用的。还有宗徒的一项特别教导与此相符：即我们每人都是一个双重的人（\BibleRef{格后 4:16}）：一个是外在可见的，其自然命运是朽坏；另一个只能在内心的隐秘处感知，却能更新。如果这教导是真实的——而且它必须是真实的，因为智慧在那里发言——那么设想一种与两者完全对应的双重婚姻也并非荒谬；而且，如果有人大胆断言身体的童贞是内在婚姻的合作者和中介，这一断言也许与可能的事实相去不远。\par

% structure: p0042 source=h2 target=section reason=relative-heading-rank; base=h2

\section{第二十章}

如今，就手工劳作而言，不可能同时操练两种技艺，例如农耕和航海，或修补和木匠。若要认真从事一种，就必须放下另一种。同样，我们面前有两种婚姻可供选择：一种在肉身中成就，另一种在圣神内成就；专注其一必然导致与另一疏离。眼睛不能同时看两个目标，而必须将注意力集中在其中一个上；舌头不能同时发出两种不同语言的话语，例如在同一时刻说出一个希伯来词和一个希腊词；耳朵不能在同一时间听取事实叙述和劝勉讲论；若每个特定的声音被分别听到，它就会将自己的观念印在听者心中；但如果它们混合在一起涌入耳中，就会产生无法解开的观念混乱，一个含义相互消失在另一个之中；同理，我们的情感能力不具备那种本性，可以同时追求感官快乐和向往圣神结合；此外，这两种目的也不能通过相同的生活方式达到：节制、克治情欲、蔑视肉体的需要，是促成那一种结合的因素；而一切与之相反的东西，则是促成肉身同居的因素。正如我们面前有两位主人可供选择，我们无法同时服从两位，因为\Quote{没有人能服事两个主人}，聪明人会选择对自己最有益的那位；同样，我们面前有两种婚姻可供选择，我们无法同时缔结两者，因为\Quote{没有结婚的人挂虑主的事，但结了婚的人挂虑世俗的事}（\BibleRef{格前 7:32}），我说，健全心智的目标应是不错过选择更有益的那一种；并且也不应不知道通往这目标的路，这条道路只能通过像下面这样的比较来认识。在现世婚姻中，一个不愿显得完全微不足道的人，会非常注意身体健康、合宜装饰、家道丰裕，以及避免有关其出身或家世的不光彩揭露；因为他认为这样事情最可能如他所愿。同样，追求圣神联合的人，首先会通过\Emph{心意的更新}而显示自己是一个年轻人，身上没有一丝衰老；其次，他会显示一种家世，这在家世上富有一种值得追求的财富，不是夸耀世俗的财富，而只在天上珍宝中奢侈享受。至于门第显赫，他不会夸耀那种仅凭传承的惯例甚至落到许多无价值之人身上的显赫，而是夸耀凭自己热忱和努力的成功追求所获得的显赫；只有那些被称为\Quote{光明之子}和天主的儿女，并因他们光辉事迹被称为\Quote{来自东方的贵族}的人，才能夸耀这种显赫。他不会通过身体锻炼和饮食来获得力量和健康，而是通过一切与此相反的事，在身体的软弱中完善精神的坚强。我还可以讲述在这样的婚礼中求婚者给新娘的礼物；它们不是由会朽坏的钱财购得，而是出自灵魂特有的财富。你想知道它们的名字吗？你必须听那位优秀的新娘装饰者\Person{保禄}所说，那些在一切事上自荐的人，他们的财富在于什么。他提到其中许多无价之物，并加上\Quote{在贞洁上}（\BibleRef{格后 6:6}）；此外，所有公认的圣神果实，无论来自何处，都是这婚姻的礼物。如果有人要遵循\BibleBook{}{箴言}的劝告，娶那真正的智慧为贤内助和终身伴侣——关于她，\BibleBook{}{箴言}说\Quote{爱她，她就保护你}，\Quote{尊崇她，她就拥抱你}（\BibleRef{箴 4:6}）——那么他就要以配得上这样爱情的方式预备自己，以便穿着无瑕的衣服与所有欢乐的婚宴宾客一同欢宴，而不是因没有穿上婚宴礼服而要求坐席时被赶出去。此外，显然这论证同样适用于男人和女人，促使他们走向这样的婚姻。宗徒说：\Quote{不再有男和女}（\BibleRef{迦 3:28}）；\Quote{基督是一切，并在一切之内}（\BibleRef{哥 3:11}）；因此，同样合理的是，爱慕智慧的人应紧握他所渴求的对象，即真智慧；而依附于不朽新郎的灵魂应享有她对真智慧——就是天主——的爱。我们现在已充分揭示了圣神联合的本性，以及纯洁而属天之爱的对象。\par

% structure: p0044 source=h2 target=section reason=relative-heading-rank; base=h2

\section{第二十一章}

十分清楚，没有人能接近天主的纯洁，除非他自己先成为纯洁；因此，他必须在自己和感官快乐之间筑起一道高大坚固的隔离墙，以便在他接近天主时，自己内心的纯洁不再被玷污。这道不可攻破的墙将在于完全与一切情欲运作之事隔绝。\par

快乐在种类上是单一的，我们可从专家们那里得知；正如水从一个单一源头分流到不同渠道，它通过感官的各种途径蔓延到快乐爱好者身上；因此，凡通过任何特定感觉而屈服于随之而来的快乐的人，其心已经受到了那感觉的伤害。这与从神圣口中所出的教诲相符：\Quote{那满足了眼目情欲的人，已经在自己心里受了伤害}；因为我认为主在那特定例子中指的是任何感官；因此我们完全可以延续他的话说，\Quote{那听了而想入非非的}，以及\Quote{那触摸了而起贪念的}，\Quote{那使我们任何官能降低为快乐服务的，心里已经犯了罪。}\par

为了防止这事，我们愿意将一切节制的准则应用于自己的生活，决不让心灵耽于任何隐藏快乐诱饵的事物；但特别要警惕味觉的快乐。因为这似乎是最根深蒂固的，并且可以说是所有被禁止的享受之母。饮食的快乐导致无度的放纵，给身体带来最可怕的痛苦之惩罚；因为过度的纵欲是大多数痛苦疾病的根源。为使身体获得持续的安宁，不受饱食之痛苦的搅动，我们必须决心采取更节省的饮食规律，以每次的需要（而非快乐）作为放纵的尺度和界限。如果甘甜仍然与需要的满足混合（因为饥饿能甜化一切，并且通过食欲的强烈，它给每一种可发现的需要之供应带来快乐的滋味），我们不可因随之而来的快乐而拒绝满足，也不可将这个满足作为首要目标。在一切事上，我们必须选择合宜的量，而留下那些仅仅满足感官的东西不动。\par

% structure: p0048 source=h2 target=section reason=relative-heading-rank; base=h2

\section{第二十二章}

我们看到农夫如何将糠秕从小麦中分离出来，以便各尽其用：前者用于人的食粮，后者用于焚烧和喂养牲畜。在节制之田里劳作的农夫，也要这样分辨满足与单纯的快乐，并将后者的本性丢给野兽，如宗徒所说的，\Quote{它们的结局是焚烧}，却怀着感恩之心，按实际需要取用关乎满足的部分。然而，许多人却滑向完全相反的一种过度，并且不自觉地，因过分拘谨而辛勤地破坏了自己的意图：他们让灵魂从神圣崇高之巅峰的另一侧，坠落到愚钝思想和俗务的水平，以至于他们的心思如此专注于仅涉及身体的规条，以致于不能再行走天上的自由、仰望高天；他们的唯一倾向就是折磨和苦待肉身。因此，最好也仔细思考此事，以便同样警惕任何一方：既不要因肉身的创伤而窒息心灵，也不要因无故施加的衰弱而削弱和降低能力，以致心灵只能想到身体的痛苦；并且让每个人记住那个智慧的箴言，它警告我们不要向右或向左转。我曾听说某位我相识的医生，在解释其技艺的秘密时，说我们的身体由四种元素组成，它们并非同类，而是倾向于冲突：然而热穿透了冷，湿与干也发生了同样出乎意料的结合，每一对矛盾者通过它们与中间一对的亲属关系而接触。他对他这种自然之研究的说明，加上了极其微妙的解释。这些元素中的每一个，其本质都与它的矛盾者\Emph{正相反对}；但又具有两边另两种性质，凭借与它们的关系，与它的矛盾者接触；例如，冷和热各与湿或干结合；同样，湿和干各与热或冷结合：因此，这种相同性质，当它出现在矛盾者中时，它本身就是促成这些矛盾者结合的因素。然而，解释这种从互相分离和本性排斥到通过相似性质的中介而互相联合的变化细节，与我何干，除非为了我们提及的目的？那个目的是要补充说，这位分析身体构造的作者建议尽一切可能保持这些性质之间的平衡，因为健康事实上就在于不让任何一种性质在我们内获得统治。如果他的教义有真理，那么，为了健康的持续，我们必须确保这样一种习惯，绝不要因饮食的任何不规则而导致这些构成元素的任何部分过多或不足。驾车者，如果他要驾驭的年轻马匹不能很好地协作，他不会用鞭子催促快马，勒住慢马；也不会任那匹在车辙中受惊或性情倔强的马独自奔驰，扰乱整齐的行进；而是加快第一匹的速度，控制第二匹，用鞭子抽打第三匹，直到使它们在一段笔直的行程中均匀呼吸。同样，我们的心灵握住这身体战车的缰绳；以此身份，在青年时期，当性情之热充足时，它不会设法增加那种热度；当身体已因疾病或年迈而寒冷时，它也不会增加冷却和稀释之物；对于所有这些身体性质，它将受圣经的引导，以致实际实现它：\BibleQuote{多收的也没有剩余，少收的也没有缺乏。}它将削减在任一方向的过度长度，并小心补充缺乏甚多之处。它将防止身体因任一原因而衰弱；它将训练肉身，既不要因过度的娇养而使之狂野难驯，也不要因过度的苦修而使之病弱、松弛、无力于必要的工作。这就是节制的最高目标：它不看苦待身体，而看灵魂功能的平安运作。\par

% structure: p0050 source=h2 target=section reason=relative-heading-rank; base=h2

\section{第二十三章}

如今，对于选择了这种哲学生活的人，其生活的细节、应避免之事、应从事的操练、节制的规则、全部的训练方法以及有助于达成这一伟大目标的日常规程，已在某些为爱好细节者所写的指导手册中论述过了。然而，有一种比言语教导更清晰的指引，那就是实践；而且，当我们开始一段漫长的旅程或航行时，应当请一位导师，这并非无理之言。\Quote{但是，} 宗徒却说，\Quote{圣言离你不远；}恩宠始于家中；那里是一切德行的制造厂；那里，通过不断迈向完美至极的进步，这种生活已被精炼得无比细致；在那里，无论人们沉默还是说话，都有充足的机会通过实际实践来学习这天上的公民身分。任何脱离生活典范的理论，无论装饰得多么令人赞叹，都如同没有呼吸的雕像，徒有以色彩和颜料描绘的红润外表；但正如福音所教导的，那位既实行又教导的人，才是真正活着的，拥有美丽的光彩，且是高效而活跃的。我们必须投奔他，若我们真要按照圣经的话\Quote{持守}童贞。一个想学外语的人，无法凭自己成为合格的教师；他接受专家的教导，然后才能与外国人交谈。同样，对于这种高尚的生活，它并不遵循自然的常规，而是以其进程的新颖性脱离了自然，除非一个人将自己托付给一位亲身完美实践了这种生活的人，否则他无法得到彻底的教导；事实上，在其他一切生活职业中，候选人如果从导师那里获得他所选择主题的每一部分的科学知识，就比自行研究更有可能取得成功；而这一特殊职业并非一切如此清晰，以至于我们必须将关于其中最佳路径的判断留给自己；它是一个一旦冒险踏入未知就会立即陷入危险的地方。医学这门科学一度并不存在；它是通过人们的实验而形成，并通过他们的各种观察逐渐被揭示；治疗性和有害的药物从那些尝试过它们的人的证词中得以认识，这种区别被采纳到该艺术的理论中，因此前人细致的观察成为后人的准则；现在，任何一个学习这门艺术的人都不必冒着自己的风险去确定任何药物的效力，无论是毒药还是药材；他只需向他人学习已知的事实，然后就可以成功实践。灵魂的医药——哲学，也是如此，我们从中学习如何治疗一切能触及灵魂的软弱。我们无需通过猜测和揣摩来追寻这些疗法的知识；我们有充足的手段从那位通过长期而丰富的经验获得了我们所追求之财宝的人那里学习它们。在一切事情上，青春通常是一个令人眩晕的向导；如果未经白发苍苍的人参与商议而能成就任何重要之事，那是不容易的。即使在其他一切事业中，我们也必须根据其重要性的大小，更加谨慎地防范失败；因为在这些事中，年轻人轻率的计划也导致了损失；例如，财产损失；或者迫使他们放弃世上的地位，甚至声望。但在这一伟大而崇高的志向中，所涉及的不是财产、一时的世俗荣耀或任何外在命运——智者对这些东西（无论它们安排得好坏）很少在意——在这里，鲁莽会影响灵魂本身；我们面临着可怕的危险，不是失去那些或许还能恢复的其他东西，而是毁灭我们自身，使灵魂破产。一个挥霍或失去了祖业的人，只要他还在人世，就不会绝望，或许还能通过筹划重新获得他先前的充足；但一个将自己从这一呼召中开除的人，也剥夺了自己一切回归更好境地的希望。因此，既然大多数人在年轻且理解力尚未成型时就拥抱童贞，那么他们首要的任务就是为自己寻找一位合适的向导和这条道路的主人，免得因为现时的无知而偏离直路，在无路的荒野中开辟自己的新路径。训道者说：\Quote{两个人胜于一个人，}\BibleRef{训 4:9}；但单独一人很容易被那条通往天主的道路上出没的仇敌击败；诚然，\Quote{那孤独一人跌倒时，有祸了，因为没有另一个人扶他起来。}此前，有些人在追求更严格生活的热忱中表现出敏捷的勤快；但是，仿佛在他们做出选择的时刻就已经达到了完美，他们的骄傲遭受了震惊的跌落，他们因疯狂地自欺，以为自己内心倾向的东西就是真正的美，而遭到了绊跌。智慧称之为\Quote{懒惰的人，}\BibleRef{箴 15:19} 的就是这类人，他们用\Quote{荆棘}铺满自己的\Quote{道路}；他们认为操心遵守诫命是一种道德上的损失；他们从自己心中抹去宗徒的教导，不食自己诚实赚得的饼，反而依赖他人，将他们的懒惰本身变成了一种生活艺术。这一类人中还有那些做梦者，他们相信梦中幻象胜过福音教导，并将自己的幻想称为\Quote{启示}。由此也产生了那些\Quote{溜进人家}的人；又有另一些人认为德行在于野蛮的熊样，从未认识恒忍和谦卑精神的果实。谁能列举出所有因拒绝求助于那些有圣德名声的人而可能陷入的陷阱呢？我们曾见过这样一类苦修者，他们坚持禁食直至死亡，仿佛\Quote{天主喜欢这样的祭献}\BibleRef{希 13:16}；而另一些人则冲向完全相反的另一极端，只在名义上过独身生活，其生活方式与世俗毫无区别；因为他们不仅纵情于口腹之乐，而且公开让人知道他们家里有女人；他们将这种友谊称为兄弟之爱，仿佛他们可以用神圣的词汇掩盖自己倾向于邪恶的思想。正因他们，这纯洁而神圣的童贞誓言\Quote{在外邦人中被亵渎。}\par

% structure: p0052 source=h2 target=section reason=relative-heading-rank; base=h2

\section{第二十四章}

因此，年轻人若不为自己在修道生活中自行规定未来的道路，实为有益；实际上，当世并不缺乏圣善生活的榜样供他们效法。如今，若说曾有圣德充盈并渗透我们的世界，那正是此时；经过逐步进展，它已触及完美的最高标杆；凡在日常循迹追随这些脚步的人，便能捕捉这一光环；凡追踪此先存芬芳之气息的人，必沉浸于基督自己的馨香之中。正如一支火炬点燃后，火焰传遍所有邻近的烛台，既不减弱首光，也不在被引燃的其他点上燃起不均等的光辉；同样，圣德的生活从成就者传至其周围的人；因为先知的话确实如此：与\Quote{圣洁}、\Quote{洁净}和\Quote{特选}的人同住，自己也会变成那样。若你愿知道确凿的标记，确保你得到真正的典范，从生活中描摹并不困难。若你看见一人如此立于死亡与生命之间，从两者中选取有助于默观生活的帮助，从不因死亡的麻木而麻痹遵守一切诫命的热情，也不因已摆脱世俗野心而将双脚稳立于活人之世；——这人对一切经审查为肉身生活的事物比死者更无感觉，却对德行的成就充满生机、精力和力量，这后者正是灵性生活的确凿标记——那么，以那人为你生活的准则；让他成为你虔敬之旅的启明星，如同永不落的星辰之于领航员；效法他的青年与白发：或者，更确切地说，效法他身上的老人与少年；因为即使在他衰年时期，时光也未迟钝他灵魂的活跃，而他的青年时期也并非活跃于青年知名的活动中；他在生命的两个季节显示了奇妙的对立结合，或更确切地说，是彼此特性的交换；因为他老年时，在向善的方向显示出青年的活力，而青年时，在邪恶的方向，他的情欲毫无能力。若你愿知道他辉煌青春期的情欲是什么，你将效法他那神性爱智的热忱与光辉，这爱智从童年与其一同成长，并延续至他老年。但若你无法注视他，如同视力弱的人无法注视太阳，至少请你注视那排列于他之下的圣洁人群，他们以生命的光照为这个时代树立了榜样。天主将他们安放为环绕我们的灯塔；他们中许多人在盛年时已是年轻人，并在不间断的节欲与节制中白了头；他们在时机未到之前已是理性上的老人，品德超越了年龄。他们唯一品尝的爱是智慧之爱；并非他们的自然本能与众不同；因为\BibleQuote{肉情相反圣神}（\BibleRef{迦 5:17}），但他们切实听从了那说节制\Quote{是生命树，凡持守她的，便为有福}的人；他们乘着这生命树如同小舟，驶过存在汹涌的波涛，锚定于天主旨意的港湾；在如此美好的航行后，他们令人羡慕，如今在晴朗无云的平静中安息灵魂。他们自己已安抵于美好盼望的锚下，远在波涛喧嚣之外；并为愿意追随的人，如在高塔上燃烧的烽火，放射他们生命的光辉。我们确实有一个标记，安全引导我们渡过诱惑的海洋；何必过分好奇地追问，是否有些怀有这些思想的人仍然跌倒了，并因此绝望，仿佛成就超出了你的能力？注视那成功的人，勇敢地启航，满怀信心这旅程将顺利，在圣神的和风下航行，以基督为舵手，以喜乐为桨。因为那些\Quote{下海坐船，在大水中经营事业}的人，并不会因某人遭遇海难而妨碍他们保持喜乐；他们反而以这信心护蔽心灵，从而继续前进完成成功的壮举。这确实是世上最荒谬的事：责备在需要最精细的课程中滑倒的人，却认为那些终生活在失败与错误中的人选择了更好的部分。若一次近于罪的事如此可怕，以致你认为根本不应着手这崇高目标，那么用罪作为一生的实践，从而完全不知更纯洁的道路，又是何等可怕！你如何在充满活力的生命中服从被钉者？你如何在罪中健壮地服从那死于罪者？你既未死于世俗，也不愿接受肉身的死灭，又如何服从那命令你跟随他，并亲自在身体上背负十字架作为战利品者？你如何服从保禄，当他劝勉你\Quote{将身体献上，当作生活、圣洁、悦乐天主的祭品}，而你却\Quote{与此世同化}，未因心思的更新而改变，未在\Quote{新生活}中\Quote{行走}，反而仍遵循\Quote{旧人}的常规？你如何能成为天主的司祭，虽为此职务受傅，向天主献上礼物——礼物绝不是外来的，不是从自身以外伪造的，而是真正属于你自己的，即\BibleQuote{内在的人}（\BibleRef{弗 3:16}），他必须完美无瑕，如同羔羊要求无玷无疵？你如何能向天主献上这礼物，若你不听从律法禁止不洁之人献祭？若你渴求天主向你显现，为何你不听从梅瑟，当日他命令百姓远离婚嫁的污秽，以便承受天主的显现（\BibleRef{出 19:15}）？若这一切在你眼中微不足道——与基督同钉十字架，将自己献给天主作祭品，成为至高天主的司祭，使自己堪当全能者的显现——那么，我们还能为你设想什么更高的福分呢？既然你连这些后果也轻视了？与基督同钉十字架的结果，是我们将与他同生，与他同受光荣，与他同享王权；将我们自己献给天主的结果，是我们将从人性和人的尊位升为天使的等级；因为达尼尔这样说：\BibleQuote{有万千者侍立在他面前}（\BibleRef{达 7:10}）。凡分享真正司祭职并置身于大司祭身旁的人，自己也永远为司祭，永不再留在死亡中。再说，使自己堪当看见天主，其成果也不小于此：即堪当看见天主。确实，一切盼望、一切愿望、一切福分、天主的一切恩许，以及我们相信存在于我们的感知与知识之外的一切不可言喻的喜乐——这一切的冠冕成果，正是这个。梅瑟切望看见它，许多先知和君王也渴望看见同一事物：但只有心地纯洁的人被堪当视为有福，并被称为\Quote{有福}，正是为此：\Quote{他们将看见天主}。因此，我们愿你同样与基督同钉十字架，成为站在天主面前的圣洁司祭，在一切贞洁中成为纯洁的祭品，以你自己的圣德预备天主来临；使你也能有一颗纯洁的心，在其中看见天主，依照天主和我们的救主耶稣基督的恩许，愿光荣归于他，至于无穷之世。阿们。\par

\SourceRecord{Translated by William Moore and Henry Austin Wilson.；Nicene and Post-Nicene Fathers, Second Series；Vol. 5.；Edited by Philip Schaff and Henry Wace.；Buffalo, NY: Christian Literature Publishing Co.,；1893.；<http://www.newadvent.org/fathers/2907.htm>.}
